%% =============================================================================
%% Projeto Base - Artigo de TCC com abnTeX2
%% Customizado para a UNIJUÍ
%% Template desenvolvido e mantido por: Prof. Eldair F. Dornelles
%% A versão mais atualizada deste template pode ser obtida no GitHub:
%% https://github.com/dornelles/tcc-artigo-template-latex
%% Qualquer sugestão de melhoria será bem vinda, 
%% Contato pelo e-mail: eldair.dornelles@unijui.edu.br
%% =============================================================================

\documentclass[
    12pt,               % tamanho da fonte
    oneside,            % para impressão frente e verso (use oneside para apenas frente)
    a4paper,            % tamanho do papel
    chapter=TITLE,      % títulos de capítulos em maiúsculas
    section=TITLE,      % títulos de seções em maiúsculas
    english,            % idioma adicional para hifenização
    brazil,             % o último idioma é o principal do documento
    article             % formatação como artigo
]{abntex2}

% ==============================================================================
% PACOTES ESSENCIAIS
% ==============================================================================

% --- Codificação e fontes ---
\usepackage[T1]{fontenc}        % seleção de códigos de fonte
\usepackage[utf8]{inputenc}     % codificação do documento (conversão automática dos acentos)
% \usepackage{lmodern}          % comentado: a fonte agora está definida no arquivo unijui.sty

% --- Pacotes de citações ---
\usepackage[brazilian,hyperpageref]{backref}  % páginas com as citações na bibl
\usepackage[alf,bibjustif]{abntex2cite}        % citações padrão ABNT (bibjustif = referências justificadas)

% --- Pacotes matemáticos ---
\usepackage{amsmath}
\usepackage{amssymb}

% --- Pacotes gráficos ---
\usepackage{graphicx}           % inclusão de gráficos
% NOTA: NÃO usar \usepackage{float} com abnTeX2/memoir - conflita com \newfloat
\usepackage{subfig}             % subfiguras
\usepackage{pdfpages}           % inclusão de páginas PDF

% --- Pacotes de tabelas ---
\usepackage{booktabs}           % tabelas profissionais
\usepackage{multirow}           % células que ocupam várias linhas
\usepackage{longtable}          % tabelas longas (múltiplas páginas)
\usepackage{array}              % personalização de colunas

% --- Outros pacotes úteis ---
\usepackage{microtype}          % melhorias de justificação
\usepackage{indentfirst}        % indenta o primeiro parágrafo de cada seção
\usepackage{color}              % controle de cores
\usepackage{xcolor}             % cores avançadas
\usepackage{listings}           % listagens de código-fonte
\usepackage{enumitem}           % listas personalizadas
\usepackage{totcount}           % conta o total de floats (figuras, tabelas, etc.) no documento

% --- Customizações UNIJUÍ ---
\input{unijui.sty}

% ==============================================================================
% CONFIGURAÇÕES DO PACOTE BACKREF
% ==============================================================================
% Configurações do pacote backref
% Usado sem a opção hyperpageref de backref
\renewcommand{\backrefpagesname}{Citado na(s) página(s):~}
% Texto padrão antes do número das páginas
\renewcommand{\backref}{}
% Define os textos da citação
\renewcommand*{\backrefalt}[4]{
    \ifcase #1 %
        Nenhuma citação no texto.%
    \or
        Citado na página #2.%
    \else
        Citado #1 vezes nas páginas #2.%
    \fi}%

% ==============================================================================
% CONFIGURAÇÕES DE LISTAGENS DE CÓDIGO
% ==============================================================================
\lstset{
    basicstyle=\ttfamily\small,
    numbers=left,
    numberstyle=\tiny\color{gray},
    stepnumber=1,
    numbersep=8pt,
    backgroundcolor=\color{white},
    showspaces=false,
    showstringspaces=false,
    showtabs=false,
    frame=single,
    rulecolor=\color{black},
    tabsize=2,
    captionpos=b,
    breaklines=true,
    breakatwhitespace=false,
    keywordstyle=\color{blue}\bfseries,
    commentstyle=\color{green!60!black},
    stringstyle=\color{red},
    escapeinside={\%*}{*)},
    literate=
        {á}{{\'a}}1 {é}{{\'e}}1 {í}{{\'i}}1 {ó}{{\'o}}1 {ú}{{\'u}}1
        {Á}{{\'A}}1 {É}{{\'E}}1 {Í}{{\'I}}1 {Ó}{{\'O}}1 {Ú}{{\'U}}1
        {ã}{{\~a}}1 {õ}{{\~o}}1 {Ã}{{\~A}}1 {Õ}{{\~O}}1
        {ç}{{\c{c}}}1 {Ç}{{\c{C}}}1
        {ê}{{\^e}}1 {â}{{\^a}}1 {ô}{{\^o}}1
}
\renewcommand{\lstlistingname}{Listagem}
\providecommand{\lstlistingautorefname}{Listagem}

% ==============================================================================
% MACROS DO ARTIGO
% ==============================================================================
% Marcador visível para dados ainda não definidos (remover na versão final).
\newcommand{\pendente}[1]{\textcolor{red}{[#1]}}
% Fonte padrão das tabelas e quadros.
\newcommand{\fonteautor}{\legend{Fonte: elaborado pelo autor (2026).}}

% ==============================================================================
% DADOS DO TRABALHO
% ==============================================================================

% --- Informações do documento (preencher conforme o trabalho) ---

\titulo{Estrutura em vez de escala: avaliação de \textit{harnesses} em nível de repositório para maximizar a qualidade da saída de LLMs restritos}
% \subtitulo{Subtítulo (se houver)}

\tituloemingles{Scaffolding Over Scale: Evaluating Repository-Level Harnesses to Maximize Output Quality in Constrained LLMs}
% \subtituloemingles{Subtitle (if any)}

\autor{Juan Paulo Fricke}
\local{Ijuí/RS}
\data{2026}
\orientador{\pendente{ORIENTADOR}}
\instituicao{%
    Universidade Regional do Noroeste do\\[1mm]
    Estado do Rio Grande do Sul
}

% -- Autores para o artigo --
\autoresartigo{
    Juan Paulo Fricke
    \thanks{Estudante do Curso de Ciência da Computação -- UNIJUÍ, Ijuí-RS, Brasil. E-mail: \texttt{\pendente{E-MAIL}}} 
    \and 
    \pendente{ORIENTADOR}
    \thanks{Professor Orientador, \pendente{TITULAÇÃO DO ORIENTADOR} -- UNIJUÍ, Ijuí-RS, Brasil. E-mail: \texttt{\pendente{E-MAIL DO ORIENTADOR}}}
}

% --- Informações específicas da UNIJUÍ (definidas no unijui.sty) ---
\campus{Ijuí}
\curso{Ciência da Computação}
\professordabanca{\pendente{MEMBRO DA BANCA}}
\datadabanca{\pendente{DIA}}{\pendente{MÊS}}{2026}

% --- Preâmbulo (texto da folha de rosto e aprovação) ---
% Não adicione ponto final do texto. Na folha de rosto o ponto irá aparecer automaticamente, 
% já na folha de aprovação ao invés de ponto irá aparecer uma vírgula, seguida do texto 
% "submetido à aprovação da banca examinadora composta pelos seguintes membros:"
\preambulo{Trabalho de Conclusão de Curso apresentado ao Curso de Ciência da Computação da 
            Universidade Regional do Noroeste do Estado do Rio Grande do Sul - UNIJUÍ,
            como requisito parcial para a obtenção do título de Bacharel em Ciência da Computação}


% --- Texto de apresentação ---
\textoapresentacao{
    O presente trabalho foi elaborado no formato de artigo científico, conforme o
    Regimento do Trabalho de Conclusão de Curso do Curso de Ciência da Computação da
    UNIJUÍ, e segue o modelo de artigo de TCC adotado pelo curso.
    \pendente{PERIÓDICO: se o artigo for submetido a um periódico ou evento, indicar
    aqui o nome do veículo e anexar as suas normas}. \par
    Os dados, os enunciados das tarefas, o \textit{snapshot} do \textit{harness}
    estudado e o código do orquestrador e da análise estão disponíveis publicamente no
    repositório do trabalho; o repositório-sujeito do experimento é privado.
}

% --- Resumo e Abstract ---
\resumoemportugues{%
Agentes de código baseados em LLMs operam em repositórios que carregam a própria estrutura de apoio: regras sempre ativas (\texttt{AGENTS.md}, \texttt{CLAUDE.md}), \textit{skills} de leitura sob demanda e configuração de ferramentas. Este trabalho pergunta se um \textit{harness} desse tipo compensa a escala do modelo em tarefas reais. Foi construído um \textit{benchmark} de ablação que reproduz 14 tarefas derivadas de \textit{pull requests} de uma aplicação Laravel privada, com os testes do próprio PR como oráculo oculto, e um orquestrador reprodutível em \textit{sandbox} Docker. Seis modelos (Claude Sonnet 5.5, Claude Haiku 5.5, Muse Spark 1.3, DeepSeek v4.1 flash, Beam e gpt-oss-20b local) executaram as 14 tarefas em quatro níveis cumulativos de \textit{harness} (\texttt{none}, \texttt{agents}, \texttt{skills}, \texttt{full}), com uma repetição. A qualidade do código foi medida em três camadas: análise estática, dez regras de convenção do repositório e até três juízes LLM às cegas por mudança (309 execuções, 836 revisões). Quatro modelos resolveram 95\% ou mais das tarefas já sem \textit{harness}, e nenhum nível alterou o \textit{pass@1} (McNemar exato, $p = 1{,}0$ em 70 pares por nível), enquanto o custo por execução subiu. O \textit{harness}, porém, elevou a nota geral dos juízes (+0,31 a +0,34; $p < 0{,}001$), a de convenções (+0,46 a +0,52) e a de testes (+0,60 a +0,87), com dez casos de teste a mais por execução, sem alterar requisitos, segurança ou escopo. Nesta amostra, a escala explica a taxa de acerto e o \textit{harness} explica parte da qualidade: ele muda como o código é escrito, não se a tarefa é resolvida.
}
% Instrução: Use letras minúsculas para as palavras-chave, separe-as com ponto e vírgula (;) e encerre a última com um ponto (.)
\palavraschaveemportugues{agentes de código; grandes modelos de linguagem; engenharia de contexto; benchmark; qualidade de código; llm como avaliador.}

\resumoemingles{
LLM-based coding agents increasingly operate in repositories that ship their own scaffolding: always-on rules (\texttt{AGENTS.md}, \texttt{CLAUDE.md}), on-demand skills and tool configuration. This work asks whether such a repository-level harness compensates for model scale on real tasks. An ablation benchmark was built that replays 14 tasks derived from pull requests of a private Laravel application, using each PR's own tests as a hidden oracle, together with a reproducible orchestrator that runs every cell in a Docker sandbox. Six models (Claude Sonnet 5.5, Claude Haiku 5.5, Muse Spark 1.3, DeepSeek v4.1 flash, Beam and a local gpt-oss-20b) ran the 14 tasks at four cumulative harness levels (\texttt{none}, \texttt{agents}, \texttt{skills}, \texttt{full}), with one repetition. Code quality was measured in three layers: static analysis, ten convention rules taken from the repository and up to three blind LLM judges per change (309 runs, 836 reviews). Four models solved 95\% or more of the tasks without any harness, and no level changed pass@1 detectably (exact McNemar, $p = 1.0$ over 70 pairs per level), while the cost per run rose. The harness did, however, raise the judges' overall score (+0.31 to +0.34; $p < 0.001$), the conventions score (+0.46 to +0.52) and the tests score (+0.60 to +0.87), with about ten more test cases per run, leaving requirements, security and scope unchanged. In this sample, scale explains the pass rate and the harness explains part of the quality: it changes how the code is written, not whether the task gets solved.
}
% Instruction: Use lowercase letters for keywords, separate them with semicolons (;) and end the last one with a period (.).
\palavraschaveemingles{coding agents; large language models; context engineering; benchmark; code quality; llm-as-a-judge.}




% ==============================================================================
% CONFIGURAÇÕES DE APARÊNCIA DO PDF FINAL
% ==============================================================================
\makeatletter
\hypersetup{
    pdftitle={\@title},
    pdfauthor={\@author},
    pdfsubject={\imprimirpreambulo},
    pdfcreator={LaTeX com abnTeX2},
    pdfkeywords={abnt}{latex}{abntex}{abntex2}{artigo}{tcc},
    colorlinks=true,        % false: links em caixas; true: links coloridos
    linkcolor=black,        % cor dos links internos
    citecolor=black,        % cor dos links para bibliografias
    filecolor=black,        % cor dos links de arquivos
    urlcolor=black,         % cor dos links de URLs
    bookmarksdepth=4
}
\makeatother

% ==============================================================================
% ESPAÇAMENTOS ENTRE LINHAS E PARÁGRAFOS
% ==============================================================================
% O tamanho do parágrafo é dado por:
\setlength{\parindent}{1.3cm}
% Controle do espaçamento entre um parágrafo e outro:
\setlength{\parskip}{0.2cm}  % tente também \onelineskip



% ==============================================================================
% INÍCIO DO DOCUMENTO
% ==============================================================================
\begin{document}

% Seleciona o idioma do documento
\selectlanguage{brazil}

% Retira espaço extra obsoleto entre as frases
\frenchspacing

% --- Imprimir página de créditos e marca d'água ---
% \imprimirpaginadecreditos
% \imprimirmarcadagua

% ==========================================================================
% ELEMENTOS PRÉ-TEXTUAIS
% ==========================================================================
\pretextual

% --- Capa ---
\imprimircapa

% --- Folha de rosto ---
% (o * indica que haverá a ficha bibliográfica)
\imprimirfolhaderosto

% --- Ficha catalográfica (inserir PDF gerado pela biblioteca) ---
% \begin{fichacatalografica}
%     \includepdf{fig_ficha_catalografica.pdf}
% \end{fichacatalografica}

% --- Folha da banca ---
\imprimirfolhadabanca

% --- Agradecimentos ---
\begin{agradecimentos}
    \pendente{AGRADECIMENTOS: texto a ser escrito pelo autor}
\end{agradecimentos}
\newpage

% --- Epígrafe ---
% Sem epígrafe nesta versão. Para incluir uma, descomente o bloco abaixo.
% \begin{epigrafe}
%     \inserirepigrafe{Texto da epígrafe.}{Autor}
% \end{epigrafe}
% \newpage

% --- Lista de ilustrações (só aparece se houver figuras) ---
% Na primeira compilação os contadores podem ser 0; compile 2x para resultado correto.
\ifnum\totvalue{figcount}>0
    \pdfbookmark[0]{\listfigurename}{lof}
    \listoffigures*
    \cleardoublepage
    \newpage
\fi

% --- Lista de quadros (só aparece se houver quadros) ---
\ifnum\totvalue{quacount}>0
    \pdfbookmark[0]{\listofquadrosname}{loq}
    \listofquadros*
    \cleardoublepage
    \newpage
\fi

% --- Lista de tabelas (só aparece se houver tabelas) ---
\ifnum\totvalue{tabcount}>0
    \pdfbookmark[0]{\listtablename}{lot}
    \listoftables*
    \cleardoublepage
    \newpage
\fi

% --- Lista de abreviaturas e siglas ---
\begin{siglas}
    \item[API] \textit{Application Programming Interface}
    \item[CLI] Interface de linha de comando (\textit{Command-Line Interface})
    \item[CSS] \textit{Cascading Style Sheets}
    \item[HTTP] \textit{Hypertext Transfer Protocol}
    \item[IC] Intervalo de confiança
    \item[JSON] \textit{JavaScript Object Notation}
    \item[LLM] Grande modelo de linguagem (\textit{Large Language Model})
    \item[LSP] \textit{Language Server Protocol}
    \item[PR] \textit{Pull request}
    \item[TCC] Trabalho de Conclusão de Curso
    \item[UI] Interface de usuário (\textit{User Interface})
    \item[UNIJUÍ] Universidade Regional do Noroeste do Estado do Rio Grande do Sul
\end{siglas}

% --- Lista de símbolos ---
% Não há símbolos neste trabalho.

% --- Sumário ---
\pdfbookmark[0]{\contentsname}{toc}
\tableofcontents*
\cleardoublepage


% --- Página de Apresentação ---
\begin{apresentacao}
    \imprimirapresentacao
\end{apresentacao}
\newpage

% --- Página de Resumo e Abstract ---
\imprimirpaginaresumos


% ==========================================================================
% ELEMENTOS TEXTUAIS
% ==========================================================================
\textual

% --- Inclusão dos capítulos ---
% ==============================================================================
% CAPÍTULO 1 - INTRODUÇÃO
% ==============================================================================

\section{Introdução}
\label{sec:introducao}

Agentes de código baseados em grandes modelos de linguagem (LLMs, do inglês
\textit{Large Language Models}) deixaram de completar linhas isoladas para
operar diretamente sobre repositórios: leem arquivos, editam código, executam
testes e iteram até considerar uma tarefa concluída. Ferramentas como o Claude
Code \cite{anthropic2025claudecode} e o OpenCode \cite{opencode2025docs}
materializam esse modo de trabalho em linha de comando, e \textit{benchmarks}
como o SWE-bench \cite{jimenez2024swebench} mostram que a taxa de resolução de
tarefas reais cresce com a capacidade do modelo. Essa capacidade, porém, tem
preço: os modelos de ponta custam mais por \textit{token}, dependem de cotas de
uso e nem sempre podem ser empregados, seja por restrições de custo, de
privacidade dos dados ou pela necessidade de execução local em
\textit{hardware} modesto.

Em paralelo, consolidou-se a prática de os próprios repositórios carregarem uma
estrutura de apoio para os agentes que neles trabalham. Arquivos de instrução
como \texttt{AGENTS.md} \cite{agentsmd2025} e \texttt{CLAUDE.md}
\cite{anthropic2025bestpractices} descrevem convenções, comandos e regras do
projeto; \textit{skills} e modelos de código ficam disponíveis para leitura sob
demanda; e arquivos de configuração habilitam servidores de linguagem (LSP,
\textit{Language Server Protocol}) e outras ferramentas. Neste trabalho, esse
conjunto de artefatos versionados no repositório é chamado de
\textit{harness}\footnote{O termo \textit{harness}, sem tradução consolidada em
português, designa aqui o conjunto de regras, \textit{skills} e configuração de
ferramentas que o repositório oferece ao agente. A expressão ``andaimes''
(\textit{scaffolding}) do título refere-se ao mesmo conjunto, visto como
estrutura de apoio ao modelo.}. A hipótese implícita nessa prática, resumida no
lema \textit{scaffolding over scale}, é que uma boa estrutura de apoio compensa,
ao menos em parte, a escala do modelo: um modelo menor bem orientado produziria
resultados comparáveis aos de um modelo maior sem orientação. Apesar da adoção
ampla, a evidência controlada sobre esse efeito ainda é escassa
\pendente{CITAÇÃO: estudos empíricos sobre o efeito de arquivos de instrução de
repositório (AGENTS.md, CLAUDE.md, regras do Cursor) em agentes de código}.

O problema que orienta esta pesquisa é, portanto, o seguinte: um
\textit{harness} em nível de repositório compensa a escala do modelo? A
pergunta desdobra-se em três: (i) o \textit{harness} aumenta a taxa com que
agentes resolvem tarefas reais, medida por testes que o agente não vê?
(ii) o \textit{harness} muda a qualidade do código produzido, mesmo quando a
tarefa já é resolvida? e (iii) qual camada do \textit{harness}, se alguma,
carrega o efeito observado?

O objetivo geral deste trabalho é avaliar, por meio de um experimento de
ablação, o efeito de um \textit{harness} de repositório real sobre a saída de
agentes de código de diferentes escalas, em tarefas derivadas do histórico
desse mesmo repositório. Os objetivos específicos são:

\begin{enumerate}[nosep]
    \item construir um \textit{benchmark} que reproduza \textit{pull requests}
    reais de uma aplicação Laravel, usando os testes de cada PR como oráculo
    oculto;
    \item definir níveis cumulativos de \textit{harness} a partir de um
    \textit{snapshot} congelado do repositório e aplicá-los de forma idêntica a
    todas as tarefas;
    \item executar seis modelos de escalas distintas em todas as tarefas e
    níveis, em ambiente isolado e reprodutível, registrando acerto, custo e
    comportamento;
    \item medir a qualidade do código produzido em três camadas independentes:
    análise estática, aderência às convenções do repositório e revisão por
    juízes LLM às cegas;
    \item analisar estatisticamente o efeito de cada nível, pareado por tarefa
    dentro de cada modelo, e discutir os resultados à luz da hipótese
    \textit{scaffolding over scale}.
\end{enumerate}

As contribuições do trabalho são três. A primeira é um \textit{benchmark} de
ablação sobre PRs reais com oráculo oculto, cujo repositório-sujeito é privado
e, portanto, pouco sujeito à contaminação dos dados de treinamento dos modelos.
A segunda é um orquestrador reprodutível que executa a matriz
modelo~$\times$~tarefa~$\times$~nível em \textit{sandbox} Docker, com
isolamento da árvore de trabalho, detecção de contaminação, tratamento de
limites de taxa e retomada sem duplicar células. A terceira é uma análise de
qualidade de código em três camadas, com juízes LLM de famílias distintas da do
modelo avaliado, validação cruzada entre camadas e medidas de concordância
entre juízes. Os dados resultantes compreendem 309 execuções avaliadas, de um
desenho de 336 células, e 615 revisões de juízes.

Antecipando os resultados, o experimento mostrou que quatro dos seis modelos
resolveram 95\% ou mais das tarefas já sem nenhum \textit{harness}, que nenhum
nível alterou a taxa de acerto de forma detectável e que o \textit{harness}
encareceu cada execução. Em contrapartida, o \textit{harness} elevou de forma
consistente as notas de convenções e de testes atribuídas pelos juízes e o
número de casos de teste escritos, sem alterar requisitos, segurança ou escopo.
O efeito concentrou-se nos modelos de desempenho intermediário e não apareceu
nos mais fracos.

O restante do artigo está organizado da seguinte forma. A
\autoref{sec:metodologia} classifica a pesquisa e descreve os procedimentos, as
métricas, os testes estatísticos e a mudança de planejamento ocorrida durante o
experimento. A \autoref{sec:referencial-teorico} situa o trabalho em cinco
eixos: \textit{benchmarks} de agentes de código, engenharia de contexto,
ferramentas de agentes, LLM como avaliador e métricas de qualidade de código. A
\autoref{sec:desenvolvimento} apresenta o objeto de estudo, o \textit{benchmark}
e o orquestrador, e em seguida os resultados de \textit{pass@1}, custo e
qualidade, a discussão e as ameaças à validade. A \autoref{sec:consideracoes}
retoma os objetivos, resume as contribuições e as limitações e aponta trabalhos
futuros.

% ==============================================================================
% CAPÍTULO 2 - METODOLOGIA
% ==============================================================================

\section{Metodologia}
\label{sec:metodologia}

Esta seção descreve o delineamento da pesquisa: sua classificação, os
procedimentos seguidos, as métricas e os testes estatísticos adotados, a
mudança de planejamento ocorrida durante o experimento e o papel das
ferramentas de inteligência artificial no método. A execução técnica de cada
etapa é detalhada na \autoref{sec:desenvolvimento}.

\subsection{Classificação da pesquisa}
\label{sec:classificacao}

Quanto à natureza, trata-se de pesquisa aplicada, pois busca responder a uma
pergunta prática da engenharia de software com LLMs: se e como vale a pena
investir em um \textit{harness} de repositório. Quanto à abordagem, é
quantitativa: todas as variáveis de resposta são medidas numericamente e
analisadas com estatística descritiva e inferencial. Quanto aos objetivos, é
explicativa, no sentido de \citeonline{gil2017}, pois procura identificar o
efeito de um fator manipulado sobre variáveis de resposta. Quanto aos
procedimentos, é um experimento controlado, na terminologia de
\citeonline{wohlin2012}, organizado como um estudo de ablação: a partir de um
\textit{harness} completo, camadas são removidas progressivamente e o efeito de
cada remoção é medido.

O fator manipulado é o nível do \textit{harness}, com quatro valores
cumulativos (\texttt{none}, \texttt{agents}, \texttt{skills} e \texttt{full},
descritos na \autoref{sec:objeto}). O modelo de linguagem é um fator de bloco:
as comparações entre níveis são feitas dentro de cada modelo, nunca entre
modelos, porque os modelos foram executados por \textit{runners} diferentes. A
tarefa é a unidade de pareamento: cada tarefa é executada por cada modelo em
cada nível, de modo que o efeito de um nível é estimado pela diferença, tarefa a
tarefa, em relação ao nível \texttt{none}. Uma célula do experimento é,
portanto, a tripla modelo~$\times$~tarefa~$\times$~nível, e uma execução
(\textit{run}) é a materialização de uma célula.

\subsection{Procedimentos}
\label{sec:procedimentos}

O experimento seguiu oito etapas, executadas entre 26 e 29 de outubro de 2026.

\textbf{Seleção das tarefas a partir de PRs.} O repositório-sujeito é uma
aplicação web de gestão de filas de atendimento escrita em Laravel (PHP) com
Bulma CSS e Livewire, privada e desenvolvida pelo autor. Cada tarefa reproduz
um \textit{pull request} (PR) ou \textit{branch} desse repositório: a árvore
no \textit{commit} base do PR é o ponto de partida do agente, um enunciado
escrito a partir do PR e da documentação de domínio é a tarefa, e os arquivos
de teste adicionados ou modificados pelo PR são o oráculo oculto, copiados para
a árvore apenas no momento da avaliação. O enunciado carrega o contrato público
de que os testes dependem (nomes de classes, métodos, rotas, chaves de
configuração), para que a tarefa meça engenharia e não adivinhação de nomes;
os pontos em que o acoplamento permanece forte são registrados por tarefa.
Dois PRs grandes, sem \textit{commit} intermediário aproveitável, foram
mantidos como tarefas compostas de referência e divididos em subtarefas que
compartilham base e solução, mas avaliam subconjuntos separáveis dos testes.
O conjunto final tem 14 tarefas: nove PRs ou \textit{branches} inteiros e
cinco subtarefas das duas compostas (\autoref{sec:benchmark} e Apêndice~A).

\textbf{Congelamento do \textit{harness}.} O \textit{harness} estudado foi
extraído de uma \textit{tag} anotada do repositório-sujeito
(\texttt{tcc-baseline-v1}, de 26 de outubro de 2026) e copiado para o
repositório do \textit{benchmark} em 27 de outubro, byte a byte. Todas as
tarefas, independentemente do ponto da história de onde vêm, recebem esse mesmo
\textit{snapshot}, o que elimina a variação do \textit{harness} entre tarefas e
isola o fator estudado.

\textbf{Definição dos níveis.} Os quatro níveis são cumulativos: \texttt{none}
não contém nenhum artefato do \textit{harness}; \texttt{agents} adiciona as
regras sempre ativas (\texttt{AGENTS.md} e \texttt{CLAUDE.md}); \texttt{skills}
adiciona o diretório \texttt{ai/} com \textit{skills}, comandos, regras e
modelos de código para leitura sob demanda; \texttt{full} adiciona a
configuração dos \textit{runners} e dos servidores de linguagem
(\texttt{opencode.json} e \texttt{.claude/settings.json}). O \textit{harness} é
aplicado em um \textit{commit} separado, marcado com a \textit{tag}
\texttt{bench/start}, de modo que o \textit{diff} do agente nunca inclui os
arquivos do \textit{harness}.

\textbf{Modelos e \textit{runners}.} Foram avaliados seis modelos de escalas e
origens distintas. Claude Sonnet 5.5 e Claude Haiku 5.5 foram executados
exclusivamente pelo Claude Code 2.1.312, com esforço de raciocínio
\texttt{high}. Muse Spark 1.3, DeepSeek v4.1 flash e Beam (plano gratuito) foram
executados pelo OpenCode v2.2.0 por meio do provedor OpenCode Go, com os
parâmetros padrão do provedor. O gpt-oss-20b \cite{openai2025gptoss} foi
executado localmente, pelo OpenCode contra um servidor llama.cpp
\cite{llamacpp2025} em uma GPU de 16~GB, com quantização MXFP4, contexto de
128 mil \textit{tokens} e a amostragem recomendada pelo fabricante. Os modelos
da Anthropic rodaram apenas no Claude Code, nunca por intermédio do OpenCode,
para que o \textit{runner} nativo de cada família fosse usado. Em cada
execução, as mensagens da sessão foram conferidas para garantir que somente o
modelo sob teste respondeu.

\textbf{Execução orquestrada e avaliação automática.} Um orquestrador escrito
em Python para este trabalho prepara cada célula, executa o agente em um
contêiner Docker com acesso apenas à própria árvore e avalia o resultado em um
segundo contêiner, copiando os testes ocultos e executando o Pint, a suíte
completa e os testes ocultos isoladamente. Todos os agentes recebem o mesmo
\textit{prompt}: um preâmbulo não interativo fixo seguido do enunciado sem
alterações, com 30 minutos de limite; a ordem das células foi embaralhada com
semente fixa (20261028) e cada execução registra versões, imagens e o
\textit{hash} do \textit{prompt}. A \autoref{sec:orquestrador} detalha o
orquestrador.

\textbf{Análise de qualidade em três camadas.} Como a taxa de acerto ficou
próxima do teto para a maioria dos modelos, uma segunda fase perguntou se o
\textit{harness} muda \emph{como} o código é escrito. Essa fase trabalha sobre
os \textit{patches} e as árvores deixadas pelas execuções, sem rodar agentes de
novo, em três camadas: (1) métricas estáticas, com Larastan nível 5
\cite{larastan2025} para contar erros novos em relação à base, complexidade
ciclomática de McCabe por função, casos de teste escritos e distância da
solução de referência; (2) dez regras heurísticas (C1 a C10) derivadas do
próprio \texttt{AGENTS.md}, cada uma aplicada apenas quando a mudança toca o
seu assunto; (3) revisão às cegas por até três juízes LLM por mudança, Claude
Opus 5.5, Muse Spark 1.3 e DeepSeek v4.1 flash, seguindo uma rubrica de sete
dimensões com notas de 1 a 5 (Apêndice~B). Cada mudança recebeu dois juízes de
famílias distintas da do modelo que a escreveu e, quando existe, o juiz da
própria família; a nota da execução é a média de todos os seus juízes, e a
nota do juiz da família é comparada à dos dois de fora como controle de
autopreferência. O \autoref{qdr:camadas} resume as três camadas. Um piloto com
duas tarefas (46 execuções, 92 revisões) validou o método antes da rodada
completa (309 execuções, 836 revisões), e todas as
violações da camada 2 apontadas no piloto foram lidas manualmente contra o
código.

\begin{quadro}[htb]
    \centering
    \ABNTEXfontereduzida
    \caption{\label{qdr:camadas}As três camadas da análise de qualidade de código}
    \begin{tabular}{|l|p{4.6cm}|p{3.4cm}|p{4.2cm}|}
        \hline
        \textbf{Camada} & \textbf{O que mede} & \textbf{Ferramenta} & \textbf{Saída por execução} \\
        \hline
        1. Estática & Erros novos de análise estática em relação à base; complexidade ciclomática por função; casos de teste e linhas escritas, separadas em produção e testes; distância da solução de referência & Larastan nível 5, analisador de complexidade em PHP, contagem de testes Pest, índice de Jaccard dos arquivos de produção & \texttt{larastan\_new\_errors}, \texttt{cc\_mean}, \texttt{test\_cases\_added}, \texttt{ref\_files\_jaccard} \\
        \hline
        2. Convenções & Dez regras do \texttt{AGENTS.md} (C1 a C10), cada uma aplicada só quando a mudança toca o seu assunto & Heurísticas sobre o \textit{patch}, com evidência guardada por violação & Escore de 0 a 1 (fração das regras aplicáveis cumpridas) \\
        \hline
        3. Juízes & Sete dimensões (requisitos, convenções, segurança, testes, simplicidade, legibilidade, escopo) e nota geral, de 1 a 5, com evidência obrigatória & Até três juízes LLM às cegas (Opus 5.5, Muse Spark 1.3, DeepSeek v4.1 flash): dois de famílias distintas da do autor e, quando existe, o da própria família; rubrica e esquema JSON & \texttt{review.json} por juiz; média de todos os juízes por dimensão \\
        \hline
    \end{tabular}
    \fonteautor
\end{quadro}

\textbf{Análise estatística.} Descrita na \autoref{sec:estatistica}.

\subsection{Métricas}
\label{sec:metricas}

A métrica principal de acerto é o \textit{pass@1}: a célula passa quando todos
os testes ocultos passam em uma única tentativa, na definição de
\citeonline{chen2021codex} para $k = 1$. Como métrica contínua, registra-se a
fração de testes ocultos aprovados. Regressões são testes fora do conjunto
oculto que falham além dos que já falhavam na árvore intocada, medidos por
tarefa com um agente nulo como linha de base. O custo é medido por turnos do
agente, \textit{tokens} de entrada (incluindo leitura de \textit{cache}) e de
saída e tempo de parede; turnos e \textit{tokens} são comparáveis entre níveis
de um mesmo modelo, enquanto o tempo depende do provedor e da concorrência. A
qualidade é medida pelas três camadas descritas acima: escore de convenções
(fração das regras aplicáveis cumpridas), erros novos do Larastan,
complexidade média, casos de teste novos, linhas inseridas em código de
produção e em testes, índice de Jaccard entre os arquivos
de produção tocados e os da referência, e as notas dos juízes em cada dimensão
e na nota geral, com a média de todos os juízes da execução.

\subsection{Análise estatística}
\label{sec:estatistica}

Todas as comparações são pareadas por tarefa, dentro de cada modelo, contra o
nível \texttt{none}. Para o \textit{pass@1}, um resultado binário, usa-se o teste
exato de McNemar \cite{mcnemar1947}, que considera apenas as tarefas
discordantes entre os dois níveis; os intervalos de confiança de 95\% das
proporções são de Wilson \cite{wilson1927}. Para as métricas contínuas e ordinais
de qualidade, usa-se o teste exato dos postos sinalizados de Wilcoxon
\cite{wilcoxon1945}, adequado a amostras pequenas sem suposição de normalidade.
A concordância entre juízes é medida pelo $\kappa$ ponderado de
\citeonline{cohen1968}, com pesos quadráticos, interpretado pela escala de
\citeonline{landis1977}, e pela correlação de postos de Spearman
\cite{spearman1904}. A autopreferência é medida, em cada mudança, pela
diferença entre a nota do juiz da família do autor e a média dos dois juízes de
fora (Wilcoxon exato pareado), descontada a leniência geral desse juiz nas
mudanças de outras famílias. A validade das notas dos juízes é verificada pela
correlação de Spearman entre a nota e medidas independentes (acerto nos testes
ocultos, escore de convenções e erros do Larastan). O agregado entre modelos
usa os 70 pares por nível dos cinco modelos com a matriz completa (5 modelos
$\times$ 14 tarefas), excluindo o Beam, que ficou incompleto. Não foi aplicada
correção para comparações múltiplas; o número de comparações é informado e a
leitura dos $p$-valores considera esse fato (\autoref{sec:ameacas}).

\subsection{Mudança de planejamento: da rodada completa à rodada reduzida}
\label{sec:mudanca}

O plano original previa cinco modelos, 14 tarefas, quatro níveis e três
repetições (840 células), com os modelos GLM 5.3 flash e MiMo v2.6 flash no
lugar de Muse Spark, DeepSeek e Beam. A rodada foi iniciada em 28 de outubro
de 2026 às 18h51 e cancelada às 19h40, depois de 41 células, quando a
estimativa de duração com tempos reais ultrapassou 30 horas; o gargalo eram os
dois modelos mais lentos do provedor, com cerca de 75 e 45 \textit{tokens} por
segundo em sessões agênticas. Com o limite de 12 horas disponível, optou-se
por uma rodada reduzida: uma repetição em vez de três, os três modelos mais
rápidos do provedor (Muse Spark, Beam e DeepSeek, medidos em 376, 339 e 280
\textit{tokens} por segundo isoladamente) e os dois modelos locais. Essa
decisão foi registrada com data e justificativa no diário do orquestrador. A
consequência metodológica é que a repetição vem das tarefas, não de execuções
repetidas da mesma célula: a variação entre tarefas tende a superar a variação
entre execuções da mesma tarefa, mas a variância de um mesmo modelo na mesma
célula não é estimada. Durante a rodada reduzida, o segundo modelo local (Qwen
3.6 35B-A3B) foi retirado por lentidão (cerca de 23 minutos por célula, o que
levaria 20 horas) e o provedor do Beam ficou fora do ar; o Beam foi encerrado
com 29 das 56 células e é apresentado apenas de forma descritiva. O desenho
final tem 6 modelos $\times$ 14 tarefas $\times$ 4 níveis $\times$ 1 repetição
(336 células).

\subsection{Uso de inteligência artificial no método}
\label{sec:uso-ia}

A inteligência artificial participa deste trabalho de três maneiras, que
convém distinguir. Primeiro, os modelos de linguagem são o \emph{objeto} do
experimento: os seis modelos avaliados escrevem o código que é medido.
Segundo, modelos são parte do \emph{instrumento} de medida: Claude Opus 5.5,
Muse Spark 1.3 e DeepSeek v4.1 flash atuam como juízes na terceira camada da
análise de qualidade, sob as salvaguardas descritas na
\autoref{sec:qualidade-metodo}. Terceiro, modelos foram \emph{ferramentas} de
construção: o orquestrador, os \textit{scripts} de análise e os relatórios
intermediários foram desenvolvidos com o auxílio do Claude Code (Claude Opus
5.5), sob direção e revisão do autor; o código do repositório-sujeito foi
escrito, em parte, com o Muse Spark, o que é discutido como ameaça à validade;
e o texto deste artigo foi redigido com o auxílio do Claude Fable 5.1 e revisado
pelo autor, que responde integralmente pelo conteúdo. A declaração formal de uso
consta ao fim do trabalho.

% ==============================================================================
% CAPÍTULO 3 - REFERENCIAL TEÓRICO
% ==============================================================================
    
\section{Referencial Teórico} 
\label{sec:referencial-teorico}

Esta seção organiza o estado da arte em cinco eixos que sustentam o desenho do
experimento: a avaliação de agentes de código em repositórios reais, a
engenharia de contexto e os arquivos de instrução de repositório, as
ferramentas que executam esses agentes, o uso de LLMs como avaliadores e as
métricas de qualidade de código. Ao fim de cada eixo, indica-se o que este
trabalho acrescenta.

\subsection{Benchmarks de agentes de código em repositórios reais}
\label{sec:ref-benchmarks}

A avaliação de geração de código por LLMs começou com problemas isolados.
\citeonline{chen2021codex} introduziram o HumanEval e a métrica
\textit{pass@k}, a probabilidade de que ao menos uma entre $k$ amostras passe
nos testes do problema, e \citeonline{austin2021mbpp} seguiram o mesmo formato
com o MBPP. \citeonline{liu2023evalplus} mostraram, porém, que os testes
desses conjuntos eram insuficientes: ao ampliá-los no EvalPlus, a taxa de acerto
de diversos modelos caiu de forma expressiva, o que evidencia a dependência de
qualquer \textit{benchmark} da qualidade do seu oráculo.

O passo para tarefas em escala de repositório foi dado pelo SWE-bench
\cite{jimenez2024swebench}, que reúne \textit{issues} reais de projetos Python
no GitHub e avalia a correção pelos testes do PR que as resolveu, separando
testes que devem passar a falhar-para-passar dos que devem continuar passando.
A versão SWE-bench Verified \cite{openai2024swebenchverified} filtrou
manualmente os itens com enunciados subespecificados ou testes injustos,
reconhecendo que parte das falhas medidas vinha do \textit{benchmark} e não dos
modelos. Sobre essa base, \citeonline{yang2024sweagent} mostraram, com o
SWE-agent, que a interface entre o agente e o computador (comandos de busca,
edição com verificação sintática, visualização paginada de arquivos) altera
substancialmente a taxa de resolução para um mesmo modelo, e
\citeonline{wang2024openhands} consolidaram no OpenHands uma plataforma aberta
para agentes generalistas de desenvolvimento. Em paralelo, trabalhos de
completação em nível de repositório, como o de \citeonline{shrivastava2023repo}
e o RepoCoder de \citeonline{zhang2023repocoder}, demonstraram que selecionar
contexto relevante do restante do repositório melhora a geração, o que
fundamenta a ideia de que o repositório pode e deve informar o modelo.

Esses trabalhos têm em comum o uso de repositórios públicos, muitas vezes
presentes nos dados de treinamento dos modelos avaliados, e a concentração em
Python. Este trabalho acrescenta um \textit{benchmark} construído sobre um
repositório privado em PHP/Laravel, cujos \textit{commits} não são acessíveis
aos modelos, com o mesmo princípio de oráculo oculto do SWE-bench, e usa esse
\textit{benchmark} não para ordenar modelos, mas para medir, por ablação, o
efeito de uma intervenção no repositório: o \textit{harness}.

\subsection{Engenharia de contexto e arquivos de instrução de repositório}
\label{sec:ref-contexto}

Agentes de código seguem, em essência, o laço de raciocínio e ação proposto por
\citeonline{yao2023react}: o modelo alterna entre pensar, chamar uma ferramenta
e observar o resultado, acumulando tudo na janela de contexto. Como essa janela
é finita e o desempenho degrada quando a informação relevante fica dispersa ou
distante \cite{liu2024lost}, a seleção do que entra no contexto passou a ser
uma disciplina própria, chamada de engenharia de contexto
\cite{anthropic2025context}. A prática recomendada é separar o que deve estar
sempre presente (regras curtas, comandos de construção e teste, convenções) do
que deve ser lido sob demanda, quando a tarefa o exigir.

Essa separação materializou-se em arquivos versionados no repositório. O
\texttt{CLAUDE.md} do Claude Code é lido automaticamente no início de cada
sessão e deve conter o mínimo necessário para o agente operar no projeto
\cite{anthropic2025bestpractices}; o \texttt{AGENTS.md} propõe-se como formato
aberto e neutro em relação à ferramenta para o mesmo fim \cite{agentsmd2025}.
\textit{Skills} são documentos mais longos, organizados por assunto, que o
agente abre quando a tarefa toca aquele assunto; a sua existência só é anunciada
nas regras sempre ativas. A configuração de ferramentas, por fim, habilita
servidores de linguagem e outros recursos que mudam o que o agente consegue
observar sem ler arquivos inteiros.

A literatura sobre esses artefatos é, em grande parte, prescritiva: descreve o
que colocar neles, mas raramente mede o efeito em condições controladas
\pendente{CITAÇÃO: estudos empíricos sobre arquivos de instrução de repositório
e sua eficácia}. Este trabalho acrescenta uma ablação em quatro níveis
cumulativos, aplicada a um \textit{harness} real e congelado, com o mesmo
enunciado e o mesmo ambiente em todos os níveis, de modo a atribuir qualquer
diferença observada à camada adicionada.

\subsection{Ferramentas e agentes de código}
\label{sec:ref-ferramentas}

O Claude Code \cite{anthropic2025claudecode} e o OpenCode
\cite{opencode2025docs} são agentes de terminal: recebem uma tarefa em
linguagem natural, dispõem de ferramentas de leitura, busca, edição e execução
de comandos, e trabalham até concluir. Ambos leem arquivos de instrução do
repositório e podem ser executados em modo não interativo, o que os torna
adequados a um experimento automatizado. Eles diferem no \textit{prompt} de
sistema, no conjunto de ferramentas e na gestão de contexto, o que torna a
comparação entre modelos executados por \textit{runners} distintos uma
comparação de pares modelo-\textit{runner}, e não de modelos isolados.

A escolha do modelo, por sua vez, envolve uma relação entre capacidade e
custo descrita pelas leis de escala: o desempenho melhora de forma previsível
com parâmetros, dados e computação \cite{kaplan2020scaling}, e há um tamanho
ótimo de modelo para cada orçamento de treinamento \cite{hoffmann2022chinchilla}.
Para quem usa os modelos, isso se traduz em famílias com variantes de tamanhos e
preços distintos e em modelos abertos, como o gpt-oss-20b
\cite{openai2025gptoss}, que podem ser executados localmente com
\textit{runtimes} como o llama.cpp \cite{llamacpp2025}, com quantização e as
limitações de \textit{hardware} que isso implica.

Este trabalho acrescenta uma comparação do mesmo \textit{harness} em seis
modelos, de uma família comercial de ponta a um modelo aberto quantizado em
GPU de consumo, executados por dois \textit{runners} sob um mesmo protocolo,
com a comparação de níveis sempre feita dentro de cada par modelo-\textit{runner}.

\subsection{LLM como avaliador}
\label{sec:ref-juiz}

Quando a métrica de interesse não é binária, como a qualidade de um
\textit{pull request}, a avaliação humana é cara e a avaliação automática
tradicional é limitada. \citeonline{zheng2023judging} propuseram usar LLMs
fortes como juízes e mostraram, com o MT-Bench e a Chatbot Arena, que o GPT-4
concorda com avaliadores humanos em nível comparável ao da concordância entre
humanos. O mesmo trabalho catalogou vieses do método: viés de posição (a
preferência pela resposta apresentada primeiro), viés de verbosidade (a
preferência por respostas mais longas) e autopreferência (a preferência por
respostas do próprio modelo). \citeonline{wang2024fair} mostraram que a simples
troca de ordem das respostas pode inverter o veredito, e
\citeonline{panickssery2024self} demonstraram que a autopreferência está ligada
à capacidade do modelo de reconhecer os próprios textos. Para o viés de
verbosidade, \citeonline{dubois2024length} propuseram controlar o comprimento
na análise. No domínio de código, \citeonline{zhuo2024icescore} mostraram que
LLMs instruídos com critérios explícitos avaliam utilidade e correção funcional
com correlação superior à de métricas baseadas em similaridade textual.

Este trabalho aplica essas lições ao desenho dos juízes: usa dois juízes por
mudança, nunca da mesma família do modelo que escreveu o código, para mitigar a
autopreferência; avalia cada mudança isoladamente, com nota absoluta e
evidência obrigatória por dimensão, para evitar o viés de posição; instrui
explicitamente a não recompensar o comprimento; e, sobretudo, mede a
concordância entre juízes e a correlação das notas com medidas independentes
(testes ocultos, regras de convenção e análise estática) em vez de assumir a
sua validade. A ligação entre um juiz (Muse Spark) e o repositório-sujeito é
discutida como ameaça à validade.

\subsection{Métricas de qualidade de código}
\label{sec:ref-qualidade}

Qualidade de código é um construto multidimensional. A norma ISO/IEC~25010
\cite{iso25010} organiza a qualidade de produto em características como
adequação funcional, segurança e manutenibilidade, e esta última desdobra-se
em analisabilidade, modificabilidade e testabilidade. Medidas estruturais
tentam aproximar essas características: a complexidade ciclomática de
\citeonline{mccabe1976} conta os caminhos linearmente independentes de uma
função e é usada há décadas como indicador de dificuldade de teste e de
manutenção. A análise estática complementa essas medidas ao apontar erros de
tipo, chamadas a métodos inexistentes e outros defeitos sem executar o código;
no ecossistema PHP, o PHPStan e sua extensão para Laravel, o Larastan
\cite{larastan2025}, cumprem esse papel em níveis de rigor crescentes.

Para código gerado por LLMs, a literatura tem apontado que correção funcional
não implica qualidade. \citeonline{pearce2022asleep} encontraram
vulnerabilidades em cerca de 40\% dos cenários em que o GitHub Copilot produziu
código relevante para segurança, o que motiva uma dimensão de segurança
separada da de requisitos. Práticas de engenharia como o desenvolvimento
orientado a testes \cite{beck2002tdd} e os princípios de legibilidade e
simplicidade de \citeonline{martin2008clean} fornecem o vocabulário das
convenções que os repositórios codificam em seus arquivos de instrução e que
revisores humanos cobram em \textit{pull requests}.

Este trabalho acrescenta uma operacionalização em três camadas que combina
medidas objetivas (Larastan, complexidade, testes escritos), regras
heurísticas derivadas das convenções do próprio repositório e julgamento por
LLM em sete dimensões alinhadas à ISO/IEC~25010 e às convenções do projeto,
com verificação cruzada entre as camadas. A escolha de medir o que o
\textit{harness} prescreve (testes, convenções, autorização, simplicidade) é
deliberada: é nessas dimensões que uma intervenção textual tem como agir.

% ==============================================================================
% CAPÍTULO 4 - DESENVOLVIMENTO
% ==============================================================================

\section{Desenvolvimento}
\label{sec:desenvolvimento}

Esta seção apresenta o objeto de estudo, o \textit{benchmark} e o orquestrador
construídos para avaliá-lo e, em seguida, os resultados de \textit{pass@1} e
custo, os resultados de qualidade de código, a discussão e as ameaças à
validade.

%-------------------------------------------------------------------------------
\subsection{O objeto de estudo: o \textit{harness} v1}
\label{sec:objeto}

O \textit{harness} estudado é o conjunto de artefatos que o repositório-sujeito
oferece aos agentes de código, congelado na \textit{tag}
\texttt{tcc-baseline-v1}. Ele tem três camadas, que correspondem aos três níveis
de tratamento do experimento (\autoref{qdr:niveis}).

A primeira camada são as regras sempre ativas. O arquivo \texttt{AGENTS.md}, com
cerca de 10~KB, declara-se autossuficiente e reúne tudo o que um agente deve
seguir no projeto: a disciplina de ``escrever menos, corretamente'' (reutilizar
antes de criar, preferir recursos nativos, nunca cortar validação ou
segurança); o desenvolvimento orientado a testes como obrigação
(``nenhum código de produção sem teste''), com execução dos testes por um
comando fixo; convenções Laravel (geradores do \texttt{artisan}, migrações
geradas e reversíveis, validação com \texttt{\#[Validate]} em componentes
Livewire, autorização por \textit{policies}, escritas de domínio em classes de
ação sob \texttt{app/Actions}, nenhuma camada de serviço ou repositório,
\texttt{\$fillable} ou \texttt{\$guarded} definidos de propósito); convenções
de interface com Bulma (CSS puro, sem bibliotecas de componentes em
JavaScript, componentes Blade com escape automático); uma política de
segurança para Livewire que trata toda propriedade pública como campo
controlável pelo cliente e exige autorização em toda ação; regras para
integrações de pagamento e para pontos de injeção de \textit{prompt}; e a
obrigação de manter a documentação e o \texttt{CHANGELOG.md} em sincronia. O
arquivo \texttt{CLAUDE.md} contém uma única linha que importa o
\texttt{AGENTS.md}, para que o Claude Code leia as mesmas regras.

A segunda camada é o diretório \texttt{ai/}, de leitura sob demanda. Ele
contém quatro \textit{skills} (\texttt{laravel-generate}, \texttt{livewire} com
a política de segurança completa, \texttt{bulma-ui} e \texttt{asaas}), comandos
equivalentes para outras ferramentas, regras no formato do Cursor e modelos de
código para controladores, componentes Livewire, modelos, \textit{requests},
testes e marcação Bulma. O \texttt{AGENTS.md} anuncia esses arquivos em uma
tabela que indica quando cada um deve ser lido; nenhuma ferramenta os carrega
automaticamente.

A terceira camada é a configuração de ferramentas. O \texttt{opencode.json}
lista o \texttt{AGENTS.md} como instrução e configura dois servidores de
linguagem, o \texttt{intelephense} e o \texttt{laravel-lsp}, para arquivos PHP
e Blade; o \texttt{.claude/settings.json} habilita o \textit{plugin} de LSP
para PHP no Claude Code. A intenção declarada no \texttt{AGENTS.md} é reduzir o
consumo de \textit{tokens}: ``prefira o LSP a leituras completas'' de arquivos.

\begin{quadro}[htb]
    \centering
    \ABNTEXfontereduzida
    \caption{\label{qdr:niveis}Níveis cumulativos do \textit{harness} e artefatos presentes na árvore}
    \begin{tabular}{|l|p{10.5cm}|}
        \hline
        \textbf{Nível} & \textbf{Artefatos presentes} \\
        \hline
        \texttt{none} & Nenhum artefato do \textit{harness} (controle). \\
        \hline
        \texttt{agents} & \texttt{AGENTS.md} e \texttt{CLAUDE.md}: regras sempre ativas. \\
        \hline
        \texttt{skills} & \texttt{agents} mais \texttt{ai/}: \textit{skills}, comandos, regras e modelos de código para leitura sob demanda. \\
        \hline
        \texttt{full} & \texttt{skills} mais \texttt{opencode.json} (instruções e LSP) e \texttt{.claude/settings.json} (LSP para PHP no Claude Code). \\
        \hline
    \end{tabular}
    \fonteautor
\end{quadro}

Dois detalhes da aplicação do nível \texttt{full} merecem registro. O
\texttt{opencode.json} original aponta o modelo padrão para o Muse Spark; o
orquestrador sobrescreve esse campo com o modelo sob teste e confere, nos
eventos da sessão, que nenhum outro modelo respondeu. O \textit{plugin} de LSP
do Claude Code só carrega quando instalado no diretório do usuário; o
orquestrador o instala em todas as execuções, de modo que o ambiente é
idêntico entre níveis e o \textit{plugin} só é ativado onde a árvore o habilita.
No OpenCode, o \texttt{laravel-lsp} foi instalado no \textit{sandbox} na mesma
versão usada pelo repositório-sujeito.

%-------------------------------------------------------------------------------
\subsection{O \textit{benchmark}: tarefas, oráculo e avaliação}
\label{sec:benchmark}

As 14 tarefas (Apêndice~A) cobrem seis categorias da aplicação: operações
(alertas no Discord e canais de \textit{log}), autenticação, multilocação
(entidades, associações e convites), domínio (filas, emissão e chamada de
senhas), interface Livewire e API. Pelo número de arquivos de produção da
solução de referência, cinco são pequenas (até 10 arquivos), sete são médias
(11 a 25) e duas são grandes (mais de 25). Duas tarefas vêm de \textit{branches}
não mescladas e são consideradas mantidas fora (\textit{held-out}). As
soluções de referência variam de 6 a 38 arquivos alterados e de 103 a 2~013
linhas inseridas. Cinco das 14 tarefas são subtarefas de dois PRs compostos:
T03a e T03b compartilham o PR de autenticação, e T04a, T04b e T04c compartilham
o PR de entidades, sendo que T04b e T04c incluem os testes de T04a. Essas
subtarefas não são observações independentes, o que é considerado nas ameaças
à validade.

Cada enunciado tem três partes: o contexto do produto, as regras de negócio e
um contrato com os nomes de que os testes ocultos dependem. O enunciado de T06,
por exemplo, especifica a sequência de códigos de senha (\texttt{A001},
\texttt{A002}, com crescimento além de 999 e tratamento de códigos legados
malformados), a rejeição de filas inativas com releitura do estado no banco
dentro da transação, a imutabilidade dos campos após a emissão e a regra de
chamada única, e lista as classes, os \textit{enums}, as \textit{factories} e as
exceções esperadas. Os enunciados não mencionam o \textit{harness} nem dizem ao
agente que ele será avaliado por testes.

O oráculo é oculto: os arquivos de teste que o PR adicionou ou modificou são
copiados da solução de referência apenas na avaliação, sobrescrevendo qualquer
arquivo de mesmo nome que o agente tenha escrito. Em cinco tarefas (T02, T05, T08, T09 e T11) a solução de
referência precisou alterar testes preexistentes, por exemplo, ao mudar o
formato dos códigos de senha; esses arquivos fazem parte do oráculo e a sua
edição pelo agente não conta como violação, enquanto a edição de qualquer outro
teste preexistente conta. A árvore entregue ao agente é um repositório Git
novo, com um único \textit{commit} de base, de modo que a solução não é
alcançável pela história, por outras \textit{branches} ou por \textit{tags}.

A avaliação, em ordem, adiciona tudo o que o agente deixou (para que arquivos
não rastreados contem), registra o \textit{diff} e as violações, copia os
testes ocultos, compila os recursos de \textit{front-end} em um contêiner Node
e, na imagem de avaliação, instala as dependências com Composer, executa o Pint,
a suíte completa e os testes ocultos isoladamente, com SQLite em memória. Se o
agente alterou o \texttt{composer.json}, as dependências são atualizadas em vez
de instaladas, o que exige acesso à rede. Um agente nulo (que não altera nada)
e um agente oráculo (que aplica a solução de referência) validaram o
avaliador: o oráculo passou em 14 de 14 tarefas e o nulo em 0 de 14. A suíte
completa tem dois testes que já falham na árvore exportada, por depender de um
diretório excluído da exportação, e essa linha de base é descontada das
regressões.

%-------------------------------------------------------------------------------
\subsection{O orquestrador}
\label{sec:orquestrador}

O orquestrador executa a matriz modelo~$\times$~tarefa~$\times$~nível em
três estágios por célula: preparação da árvore e do \textit{prompt}, execução
do agente em um \textit{pool} por provedor e avaliação em um número fixo de
\textit{slots} de Docker. Ele foi escrito em Python, com a biblioteca padrão, e
seu estado é um manifesto em linhas JSON, com uma linha por transição de
célula, que também é a fonte da retomada: a mesma invocação, repetida, pula as
células concluídas e continua as pendentes sem duplicar nenhuma.

\textbf{Isolamento.} Durante o desenvolvimento verificou-se que, sem
\textit{sandbox}, um agente com permissões automáticas alcançaria, a partir da
árvore de trabalho, os enunciados com os arquivos da solução, os
\textit{patches} de outras execuções e o próprio repositório-sujeito. Por
isso, cada agente roda em um contêiner Docker com apenas a sua árvore montada e
nenhum outro caminho do \textit{host}, com acesso completo à internet e, a
partir das 01h53 do dia 29, com limite de 6~GB de memória, introduzido depois
que um contêiner levou o \textit{host} ao limite de memória e foi morto pelo
sistema. O comando do projeto para executar testes dentro de contêineres é
substituído por um \textit{shim} que executa o comando no próprio
\textit{sandbox}, que contém PHP e Composer; assim, o agente pode rodar a
suíte como o \texttt{AGENTS.md} instrui, em todos os níveis. As credenciais
dos provedores nunca ficam no repositório: são exportadas para um diretório
temporário por execução e apagadas ao fim.

\textbf{Detecção de contaminação.} Depois de cada execução, as entradas das
chamadas de ferramenta são varridas em busca de caminhos fora da árvore e de
referências ao repositório-sujeito. Execuções suspeitas são excluídas do
\textit{pass@1} e mantidas em disco para auditoria. Durante a rodada, todos os
sinais do detector revelaram-se falsos positivos (um \textit{link} para um PR
escrito no \texttt{CHANGELOG.md} e uma navegação relativa dentro de
\texttt{vendor/}) e as regras foram corrigidas; o veredito passou a ser
recalculado a partir do \textit{log} bruto completo. Tentativas de acessar o
repositório-sujeito no GitHub passaram a ser contadas separadamente, porque o
repositório é privado e o \textit{sandbox} não tem credencial, de modo que
elas não podem vazar nada.

\textbf{Limites de taxa e falhas.} Uma resposta de limite de taxa (HTTP 429,
cota ou janela de uso) não é dado e não consome tentativa: o \textit{pool} do
provedor pausa por 60 segundos, dobrando a cada repetição até 30 minutos, e a
célula espera. Outras falhas de infraestrutura (rede, erro do servidor,
\textit{crash} da CLI, modelo errado respondendo) são refeitas com recuo e
limite de tentativas. Um \textit{timeout} de 30 minutos é dado, e a árvore é
avaliada como ficou. Uma falha na avaliação leva a uma nova avaliação da
mesma árvore, sem rodar o agente de novo. Durante a rodada foi acrescentado um
\textit{watchdog} de silêncio: um agente sem nenhum evento por oito minutos é
interrompido e tratado como limite de taxa, o que resolveu os travamentos do
provedor do Beam sem convertê-los em dados. A saída malformada do modelo local,
que o servidor rejeita, foi classificada como falha do modelo, e não de
infraestrutura, para não enviesar o gpt-oss para cima ao repetir apenas as
tentativas em que ele acertou o formato.

\textbf{Reprodutibilidade.} Cada invocação grava um instantâneo do ambiente
com as versões do OpenCode, do Claude Code, do Docker e do llama.cpp, os
identificadores das imagens, o \textit{commit} do código do orquestrador, o
\textit{hash} da matriz e a semente. Cada execução grava o \textit{hash} do
preâmbulo, o texto exato do \textit{prompt} enviado, a versão da CLI, o
esforço, o \textit{sandbox} e os modelos vistos na sessão. Testes de aceitação
verificaram que a execução sequencial e a paralela produzem o mesmo resultado,
que a retomada não duplica células e que o 429 pausa o \textit{pool}.

%-------------------------------------------------------------------------------
\subsection{Resultados: \textit{pass@1} e custo}
\label{sec:resultados-pass}

A rodada reduzida foi executada em 29 de outubro de 2026, das 00h19 às 05h30,
com o Beam retomado das 07h08 às 07h45 e encerrado parcial. A
\autoref{tab:modelos} resume os modelos e parâmetros, e a \autoref{tab:pass1}
apresenta o \textit{pass@1} por modelo e nível, com intervalos de confiança de
Wilson a 95\%.

\begin{table}[htb]
    \centering
    \ABNTEXfontereduzida
    \caption{Modelos avaliados, \textit{runners} e parâmetros}
    \label{tab:modelos}
    \begin{tabular}{llll}
        \toprule
        \textbf{Modelo} & \textbf{Runner} & \textbf{Parâmetros} & \textbf{Concorrência} \\
        \midrule
        Claude Sonnet 5.5 & Claude Code 2.1.312 & esforço \texttt{high} & 3 \\
        Claude Haiku 5.5 & Claude Code 2.1.312 & esforço \texttt{high} & 3 \\
        Muse Spark 1.3 & OpenCode v2.2.0 (Go) & padrão do provedor & 3 \\
        DeepSeek v4.1 flash & OpenCode v2.2.0 (Go) & padrão do provedor & 3 \\
        Beam (free) & OpenCode v2.2.0 (Go/Zen) & padrão do provedor & 3 \\
        gpt-oss-20b (local) & OpenCode v2.2.0 + llama.cpp & MXFP4, 128k, temp. 1,0 & 1 \\
        \bottomrule
    \end{tabular}
    \fonteautor
\end{table}

\begin{table}[htb]
    \centering
    \ABNTEXfontereduzida
    \caption{\textit{Pass@1} por modelo e nível de \textit{harness} (tarefas resolvidas em 14; IC 95\% de Wilson entre colchetes)}
    \label{tab:pass1}
    \begin{tabular}{lccccc}
        \toprule
        \textbf{Modelo} & \texttt{none} & \texttt{agents} & \texttt{skills} & \texttt{full} & \textbf{Total} \\
        \midrule
        Sonnet 5.5 & 14 [0,78; 1,00] & 14 [0,78; 1,00] & 14 [0,78; 1,00] & 13 [0,69; 0,99] & 55/56 (0,98) \\
        DeepSeek v4.1 flash & 14 [0,78; 1,00] & 14 [0,78; 1,00] & 14 [0,78; 1,00] & 14 [0,78; 1,00] & 56/56 (1,00) \\
        Muse Spark 1.3 & 13 [0,69; 0,99] & 13 [0,69; 0,99] & 14 [0,78; 1,00] & 13 [0,69; 0,99] & 53/56 (0,95) \\
        Haiku 5.5 & 11 [0,52; 0,92] & 11 [0,52; 0,92] & 11 [0,52; 0,92] & 11 [0,52; 0,92] & 44/56 (0,79) \\
        gpt-oss-20b & 2 [0,04; 0,40] & 1 [0,01; 0,31] & 1 [0,01; 0,31] & 2 [0,04; 0,40] & 6/56 (0,11) \\
        Beam (free)$^{a}$ & 10/10 & 7/7 & 5/5 & 7/7 & 29/29 (1,00) \\
        \bottomrule
    \end{tabular}
    \legend{Fonte: elaborado pelo autor (2026). $^{a}$ Beam incompleto: 29 das 56 células (23 pelo OpenCode Go e 6 pelo OpenCode Zen); as células concluídas não são uma amostra aleatória das tarefas e ficam fora das comparações entre níveis.}
\end{table}

\textbf{Efeito de teto.} Quatro dos seis modelos resolveram 95\% ou mais das
tarefas já no nível \texttt{none}: Sonnet 5.5 (55/56 no total), DeepSeek v4.1
flash (56/56), Muse Spark 1.3 (53/56) e Beam (29/29 nas células que rodaram).
Com o teto tão próximo, não há espaço para o \textit{harness} melhorar o
\textit{pass@1} desses modelos. Mesmo quando falham, os modelos fortes falham
por um único teste: a fração média de testes ocultos aprovados fica em 0,998 ou
mais para eles e em 0,993 para o Haiku 5.5 em todos os níveis.

\textbf{Nenhum efeito do \textit{harness} sobre o \textit{pass@1}.} A
\autoref{tab:pareado} mostra, para cada modelo e nível, o número de tarefas em
que apenas o nível passou ($+$) e em que apenas o \texttt{none} passou ($-$). Nos
70 pares modelo~$\times$~tarefa de cada nível (5 modelos completos $\times$ 14
tarefas), houve no máximo uma discordância para cada lado. As diferenças
agregadas foram de $-1{,}4$ pontos percentuais para \texttt{agents}, $0{,}0$ para
\texttt{skills} e $-1{,}4$ para \texttt{full}, com $p = 1{,}0$ no teste exato de
McNemar nos três casos. Em 5 modelos $\times$ 3 níveis houve apenas quatro
discordâncias, nas duas direções. O padrão da grade tarefa~$\times$~nível é
dominado pelo modelo e pela tarefa: as falhas se repetem na mesma tarefa em
todos os níveis.

\begin{table}[htb]
    \centering
    \ABNTEXfontereduzida
    \caption{Efeito de cada nível contra \texttt{none} no \textit{pass@1}, pareado por tarefa ($+$: só o nível passou; $-$: só o \texttt{none} passou)}
    \label{tab:pareado}
    \begin{tabular}{lccc}
        \toprule
        \textbf{Modelo} & \texttt{agents} ($+/-$) & \texttt{skills} ($+/-$) & \texttt{full} ($+/-$) \\
        \midrule
        Sonnet 5.5 & 0/0 & 0/0 & 0/1 (T05) \\
        DeepSeek v4.1 flash & 0/0 & 0/0 & 0/0 \\
        Muse Spark 1.3 & 0/0 & 1/0 (T05) & 0/0 \\
        Haiku 5.5 & 0/0 & 0/0 & 0/0 \\
        gpt-oss-20b & 0/1 (T07) & 0/1 (T07) & 0/0 \\
        \midrule
        \textbf{Agregado} (70 pares) & 0/1; $\Delta = -1{,}4$ pp; $p = 1{,}0$ & 1/1; $\Delta = 0{,}0$ pp; $p = 1{,}0$ & 0/1; $\Delta = -1{,}4$ pp; $p = 1{,}0$ \\
        \bottomrule
    \end{tabular}
    \fonteautor
\end{table}

\textbf{Os dois modelos abaixo do teto não reagiram ao \textit{harness}.} O
Haiku 5.5 resolveu 11 de 14 tarefas em todos os níveis e falhou sempre nas
mesmas três (T04a, T04b e T04c) pelo mesmo teste: a política de associação
deveria permitir que um administrador atribua papel ao dono da entidade, como
o enunciado diz (``\textit{owner and admin may assign roles}''), mas o modelo
acrescenta em todas as 12 células uma proteção ao dono que não foi pedida. É
uma interpretação estável do modelo, que nenhum nível de \textit{harness}
corrigiu. O gpt-oss-20b resolveu 6 de 56 células (1 ou 2 de 14 por nível), com
fração média de testes aprovados em torno de um quarto; a métrica contínua
sugere uma piora leve nos níveis \texttt{skills} e \texttt{full} (de 0,250 para
0,203), sem significância com $n = 14$. Das 56 execuções dele, 24 terminaram
com saída malformada que o servidor não conseguiu interpretar, 35 deixaram
regressões além da linha de base (em média 5,5 testes, em geral código que
quebra a aplicação inteira) e apenas 3 passaram no Pint.

\textbf{O \textit{harness} aumenta o custo.} A \autoref{tab:custo} apresenta as
medianas por execução. Em relação ao \texttt{none}, o Sonnet 5.5 passou de 18
para 26 a 30 turnos (+44\% a +67\%), de 906 mil para 1,36 a 1,51 milhão de
\textit{tokens} de entrada (+50\% a +66\%) e de 3,1 para 4,1 a 4,6 minutos de
agente (+32\% a +48\%). O Haiku 5.5 subiu 13\% a 19\% em turnos, 43\% a 71\% em
\textit{tokens} de entrada e 37\% a 53\% em tempo. O DeepSeek subiu 17\% a 62\% em
\textit{tokens} de entrada e o Muse 5\% a 21\%. O gpt-oss, que raramente abre os
arquivos do \textit{harness}, quase não mudou. Os níveis \texttt{agents},
\texttt{skills} e \texttt{full} custam de forma parecida entre si: a maior parte
do acréscimo vem da primeira camada, e as demais acrescentam pouco. Em totais,
a rodada consumiu 3,7 horas de agente e 1,57 milhão de \textit{tokens} de saída
no Sonnet (US\$~38,7 a preço de lista, segundo a estimativa da própria CLI), 7,0
horas e 4,87 milhões no Haiku (US\$~27,3), 6,1 horas no Muse, 8,1 horas no
DeepSeek e 3,0 horas no gpt-oss. O Haiku gerou cerca de três vezes mais
\textit{tokens} de saída que o Sonnet e levou cerca do dobro do tempo, com
\textit{pass@1} menor. O tempo de parede não é comparável entre modelos, pois
depende do provedor, do \textit{hardware} e da concorrência; turnos e
\textit{tokens} são a medida comparável.

\begin{table}[htb]
    \centering
    \ABNTEXfontereduzida
    \caption{Custo por execução (medianas) por modelo e nível}
    \label{tab:custo}
    \begin{tabular}{llrrrr}
        \toprule
        \textbf{Modelo} & \textbf{Nível} & \textbf{Turnos} & \textbf{Tokens de saída} & \textbf{Tokens de entrada} & \textbf{Tempo (min)} \\
        \midrule
        Sonnet 5.5 & \texttt{none} & 18 & 22\,502 & 905\,810 & 3,1 \\
         & \texttt{agents} & 30 & 30\,112 & 1\,468\,970 & 4,1 \\
         & \texttt{skills} & 29 & 31\,398 & 1\,507\,708 & 4,5 \\
         & \texttt{full} & 26 & 27\,909 & 1\,363\,034 & 4,6 \\
        \midrule
        Haiku 5.5 & \texttt{none} & 62 & 73\,469 & 3\,767\,150 & 5,9 \\
         & \texttt{agents} & 70 & 86\,608 & 5\,397\,860 & 8,3 \\
         & \texttt{skills} & 74 & 85\,562 & 6\,428\,217 & 9,0 \\
         & \texttt{full} & 71 & 87\,200 & 5\,525\,890 & 8,1 \\
        \midrule
        Muse Spark 1.3 & \texttt{none} & 62 & 20\,360 & 4\,125\,950 & 4,1 \\
         & \texttt{agents} & 64 & 18\,510 & 4\,354\,079 & 8,6 \\
         & \texttt{skills} & 74 & 18\,452 & 4\,700\,618 & 5,1 \\
         & \texttt{full} & 67 & 19\,244 & 4\,989\,247 & 5,8 \\
        \midrule
        DeepSeek v4.1 flash & \texttt{none} & 58 & 17\,518 & 4\,256\,746 & 9,3 \\
         & \texttt{agents} & 68 & 20\,720 & 4\,962\,629 & 8,7 \\
         & \texttt{skills} & 67 & 20\,570 & 6\,097\,903 & 8,3 \\
         & \texttt{full} & 76 & 22\,686 & 6\,913\,314 & 8,4 \\
        \midrule
        gpt-oss-20b & \texttt{none} & 42 & 10\,990 & 693\,686 & 2,5 \\
         & \texttt{agents} & 38 & 10\,620 & 691\,718 & 1,8 \\
         & \texttt{skills} & 46 & 13\,660 & 977\,582 & 2,4 \\
         & \texttt{full} & 46 & 11\,222 & 982\,754 & 2,3 \\
        \bottomrule
    \end{tabular}
    \legend{Fonte: elaborado pelo autor (2026). Tokens de entrada incluem leitura de \textit{cache}. O tempo é o de parede do agente e não é comparável entre modelos.}
\end{table}

\textbf{Comportamento.} Houve um único \textit{timeout} válido (DeepSeek em
T05 \texttt{none}, que trabalhou os 30 minutos e passou mesmo assim) e nenhuma
contaminação. Nas 56 execuções de cada modelo, as regressões ficaram entre 1 e
5 para os modelos de API, contra 35 para o gpt-oss; nenhum agente editou
arquivos do \textit{harness} nem introduziu saída Blade sem escape; Sonnet e
Haiku alteraram testes preexistentes em 11 e 9 execuções, DeepSeek em 3 e Muse
em nenhuma. Três execuções tentaram buscar o repositório-sujeito no GitHub
(uma leitura de arquivo bruto, uma consulta à API e um \texttt{git ls-remote}),
todas do DeepSeek e todas no nível \texttt{none}; todas falharam com erro ou
404. Com $n = 3$ o fato é anedótico, mas é coerente com a ideia de que o
\textit{harness} dá ao agente contexto suficiente para não procurar fora.

%-------------------------------------------------------------------------------
\subsection{Resultados: qualidade de código}
\label{sec:resultados-qualidade}
\label{sec:qualidade-metodo}

A análise de qualidade cobriu as 309 execuções avaliadas da rodada (56 de cada
um dos cinco modelos completos e 29 do Beam), com as três camadas aplicadas a
todas; 306 execuções receberam as duas revisões de juízes e três ficaram com
uma, porque a segunda revisão terminou em erro do juiz. Os pares de juízes
foram fixos por modelo avaliado: Muse e DeepSeek para Sonnet e Haiku; Opus e
DeepSeek para o Muse; Opus e Muse para o DeepSeek; e rotação balanceada das
três duplas para gpt-oss e Beam. Cada juiz trabalhou em um espaço isolado com
o enunciado, o \textit{diff} sem os arquivos do \textit{harness} e sem os
\textit{lock files}, as seções normativas do \texttt{AGENTS.md} como padrão de
revisão (o mesmo para todos os níveis), a rubrica e o esquema de saída. Uma
verificação recusa qualquer espaço que contenha o identificador da execução, o
modelo, o nível ou o resultado dos testes. A revisão é validada contra um
esquema JSON e, se inválida, recebe uma segunda chance com os erros
informados.

\textbf{Efeito por nível.} A \autoref{tab:qualidade-efeito} resume o efeito de
cada nível contra o \texttt{none}, pareado por tarefa, nos 70 pares dos cinco
modelos completos. O \textit{harness} elevou a nota geral dos juízes em 0,31 a
0,34 ponto (em uma escala de 1 a 5; $p \le 0{,}001$), a dimensão de convenções
em 0,44 a 0,51 ($p < 0{,}001$) e a dimensão de testes em 0,56 a 0,82
($p < 0{,}001$). Na camada estática, o número de casos de teste novos subiu 10
a 11 por execução ($p < 0{,}001$), o escore de convenções subiu 0,07 a 0,10
($p \le 0{,}001$), os erros novos do Larastan caíram 0,5 a 0,6 por execução
($p$ entre 0,015 e 0,057) e a complexidade média caiu 0,16 a 0,19 ($p$ entre
0,001 e 0,051). Requisitos, segurança e escopo não mudaram ($p$ entre 0,06 e
0,82), nem a distância da solução de referência (Jaccard, $p \ge 0{,}17$). A
legibilidade subiu pouco (0,06 a 0,13) e a simplicidade não mudou de forma
consistente. O nível \texttt{skills} tem os maiores efeitos em convenções e
testes, mas as diferenças entre \texttt{agents}, \texttt{skills} e \texttt{full}
são pequenas diante dos efeitos de cada um contra o \texttt{none}: a maior parte
do ganho vem da primeira camada.

\begin{table}[htb]
    \centering
    \ABNTEXfontereduzida
    \caption{Efeito de cada nível contra \texttt{none} nas métricas de qualidade ($\Delta$ = média das diferenças pareadas por tarefa; Wilcoxon exato; 70 pares por nível, 63 a 70 nas métricas com dados ausentes)}
    \label{tab:qualidade-efeito}
    \begin{tabular}{lrrrrrr}
        \toprule
        & \multicolumn{2}{c}{\texttt{agents}} & \multicolumn{2}{c}{\texttt{skills}} & \multicolumn{2}{c}{\texttt{full}} \\
        \cmidrule(lr){2-3} \cmidrule(lr){4-5} \cmidrule(lr){6-7}
        \textbf{Métrica} & $\Delta$ & $p$ & $\Delta$ & $p$ & $\Delta$ & $p$ \\
        \midrule
        \multicolumn{7}{l}{\textit{Camada 1 e 2 (medidas objetivas)}} \\
        Escore de convenções (0 a 1) & $+0{,}07$ & 0,001 & $+0{,}10$ & $<0{,}001$ & $+0{,}07$ & $<0{,}001$ \\
        Larastan: erros novos & $-0{,}53$ & 0,034 & $-0{,}57$ & 0,057 & $-0{,}63$ & 0,015 \\
        Complexidade média & $-0{,}17$ & 0,001 & $-0{,}19$ & 0,025 & $-0{,}16$ & 0,051 \\
        Casos de teste novos & $+10{,}80$ & $<0{,}001$ & $+11{,}43$ & $<0{,}001$ & $+10{,}14$ & $<0{,}001$ \\
        Jaccard vs. referência & $-0{,}00$ & 0,265 & $-0{,}01$ & 0,245 & $-0{,}01$ & 0,169 \\
        \midrule
        \multicolumn{7}{l}{\textit{Camada 3 (juízes, notas de 1 a 5)}} \\
        Nota geral & $+0{,}31$ & 0,001 & $+0{,}34$ & 0,001 & $+0{,}33$ & $<0{,}001$ \\
        Requisitos & $+0{,}02$ & 0,313 & $-0{,}04$ & 0,732 & $+0{,}01$ & 0,815 \\
        Convenções & $+0{,}44$ & $<0{,}001$ & $+0{,}51$ & $<0{,}001$ & $+0{,}46$ & $<0{,}001$ \\
        Segurança & $+0{,}08$ & 0,230 & $+0{,}12$ & 0,056 & $+0{,}07$ & 0,171 \\
        Testes & $+0{,}67$ & $<0{,}001$ & $+0{,}82$ & $<0{,}001$ & $+0{,}56$ & $<0{,}001$ \\
        Simplicidade & $+0{,}13$ & 0,057 & $+0{,}04$ & 0,557 & $+0{,}09$ & 0,101 \\
        Legibilidade & $+0{,}13$ & 0,020 & $+0{,}06$ & 0,195 & $+0{,}11$ & 0,042 \\
        Escopo & $+0{,}01$ & 0,659 & $-0{,}01$ & 0,638 & $+0{,}03$ & 0,559 \\
        \bottomrule
    \end{tabular}
    \legend{Fonte: elaborado pelo autor (2026). O Beam, incompleto, não entra no agregado. $p$-valores sem correção para comparações múltiplas (39 comparações).}
\end{table}

\textbf{Efeito por modelo.} A \autoref{tab:qualidade-modelo} mostra a nota
geral dos juízes e o número de casos de teste novos por modelo e nível. O
efeito é maior em quem parte mais baixo. O DeepSeek v4.1 flash é quem mais
ganha: a nota geral vai de 3,18 no \texttt{none} para 3,93 a 4,11 com
\textit{harness}, a nota de testes de 2,79 para 4,43 a 4,57 e os casos de
teste de 8,0 para 27,8 a 32,1 por execução. O Muse Spark 1.3 vem em seguida:
nota geral de 2,61 para 3,04 a 3,32, testes de 1,21 para 2,07 a 2,75, casos de
teste de 1,1 para 5,4 a 9,7. Sonnet e Haiku ganham pouco na nota geral (de
3,96 para 4,14 a 4,17 e de 4,18 para 4,25 a 4,39), porque já partem de
convenções e testes altos, mas o Sonnet dobra os casos de teste escritos (de
16,7 para cerca de 34). Os dois modelos mais fracos não respondem ao
\textit{harness}: o gpt-oss-20b fica com nota geral entre 1,43 e 1,61 e nota de
testes de 1,0 em quase todos os níveis, e o Beam, nas células disponíveis, fica
entre 2,36 e 2,64, também sem testes. Em ambos, as regras do \texttt{AGENTS.md}
que mandam escrever testes simplesmente não são seguidas.

\begin{table}[htb]
    \centering
    \ABNTEXfontereduzida
    \caption{Nota geral dos juízes (média dos dois juízes, 1 a 5) e casos de teste novos por execução, por modelo e nível}
    \label{tab:qualidade-modelo}
    \begin{tabular}{lcccccccc}
        \toprule
        & \multicolumn{4}{c}{\textbf{Nota geral dos juízes}} & \multicolumn{4}{c}{\textbf{Casos de teste novos}} \\
        \cmidrule(lr){2-5} \cmidrule(lr){6-9}
        \textbf{Modelo} & \texttt{none} & \texttt{agents} & \texttt{skills} & \texttt{full} & \texttt{none} & \texttt{agents} & \texttt{skills} & \texttt{full} \\
        \midrule
        Sonnet 5.5 & 3,96 & 4,14 & 4,17 & 4,14 & 16,7 & 33,9 & 34,3 & 33,9 \\
        Haiku 5.5 & 4,18 & 4,39 & 4,31 & 4,25 & 36,8 & 45,3 & 45,4 & 42,0 \\
        Muse Spark 1.3 & 2,61 & 3,04 & 3,32 & 3,11 & 1,1 & 9,7 & 8,7 & 5,4 \\
        DeepSeek v4.1 flash & 3,18 & 4,11 & 3,93 & 4,07 & 8,0 & 27,8 & 28,4 & 32,1 \\
        Beam (free)$^{a}$ & 2,50 & 2,36 & 2,40 & 2,64 & 0,0 & 0,0 & 0,0 & 1,1 \\
        gpt-oss-20b & 1,61 & 1,43 & 1,46 & 1,61 & 0,0 & 0,0 & 2,9 & 0,0 \\
        \bottomrule
    \end{tabular}
    \legend{Fonte: elaborado pelo autor (2026). $n = 14$ por célula, exceto Sonnet \texttt{skills} ($n = 12$) e Haiku \texttt{skills} ($n = 13$) na nota geral. $^{a}$ Beam: $n$ = 10, 7, 5 e 7.}
\end{table}

\textbf{Um exemplo.} O caso do Muse Spark na tarefa T07 (chamada determinística
da próxima senha com modos de chamada e regra de justiça) ilustra o padrão. Nos
quatro níveis o modelo passou nos testes ocultos. No nível \texttt{none}, o
\textit{diff} trouxe a ação de domínio, o \textit{enum} de modos, a exceção e
a documentação, e nenhum arquivo de teste; os dois juízes deram nota 1 em
testes e nota geral 3, com a observação de que ``a lógica crítica de
concorrência é entregue sem nenhum teste''. No nível \texttt{skills}, o mesmo
modelo entregou, além do código, um arquivo de teste com quinze casos cobrindo
ordenação e empates, escopo por fila, os três modos e seus caminhos de erro,
fila vazia, a regra de justiça (disparo, reinício, histórico de outros modos) e
fila inativa; os juízes deram nota 4 em testes e nota geral 4. O oráculo oculto
não distingue as duas entregas; a revisão de código distingue.

\textbf{Concordância e validade dos juízes.} A \autoref{tab:kappa} mostra a
concordância entre as três duplas de juízes na nota geral. O $\kappa$ ponderado
vai de 0,71 (Muse e DeepSeek) a 0,85 (Opus e Muse), o que corresponde a
concordância substancial a quase perfeita na escala de
\citeonline{landis1977}; o $\kappa$ médio nas sete dimensões vai de 0,73 a 0,84.
Os juízes têm vieses pequenos e em direções distintas: o Muse é o mais
generoso (+0,22 em relação à média da dupla), o Opus o mais severo ($-0{,}17$) e
o DeepSeek fica perto do centro ($-0{,}09$). Como o par de juízes é fixo por
modelo avaliado, o viés do juiz é constante entre os níveis de um mesmo modelo
e não afeta as comparações pareadas, mas afeta a comparação de níveis absolutos
entre modelos.

\begin{table}[htb]
    \centering
    \ABNTEXfontereduzida
    \caption{Concordância entre juízes na nota geral e nas sete dimensões}
    \label{tab:kappa}
    \begin{tabular}{lcccc}
        \toprule
        \textbf{Dupla} & $n$ & $\kappa$ ponderado (geral) & Spearman (geral) & $\kappa$ médio (7 dimensões) \\
        \midrule
        Opus $\times$ Muse & 82 & 0,85 & 0,91 & 0,73 \\
        Opus $\times$ DeepSeek & 87 & 0,81 & 0,84 & 0,84 \\
        Muse $\times$ DeepSeek & 137 & 0,71 & 0,68 & 0,79 \\
        \bottomrule
    \end{tabular}
    \fonteautor
\end{table}

A \autoref{tab:validade} verifica se as notas medem algo externo a elas. A nota
geral correlaciona-se com o acerto nos testes ocultos ($\rho = 0{,}52$), e a
dimensão de requisitos um pouco mais ($\rho = 0{,}54$): os juízes, que não veem
o resultado dos testes, tendem a dar notas mais altas às mudanças que passam.
A dimensão de convenções dos juízes correlaciona-se fortemente com o escore
heurístico da camada 2 ($\rho = 0{,}73$), duas medidas independentes do mesmo
construto. Em contrapartida, a nota geral não se correlaciona com os erros novos
do Larastan ($\rho = -0{,}07$); no piloto, com duas tarefas, essa correlação era
de $-0{,}53$, e a divergência é discutida nas ameaças à validade.

\begin{table}[htb]
    \centering
    \ABNTEXfontereduzida
    \caption{Validade: correlação de Spearman entre medidas de camadas distintas (todas as execuções)}
    \label{tab:validade}
    \begin{tabular}{llcc}
        \toprule
        \textbf{Medida A} & \textbf{Medida B} & $\rho$ & $n$ \\
        \midrule
        Juízes: nota geral & Testes ocultos aprovados & 0,52 & 306 \\
        Juízes: requisitos & Testes ocultos aprovados & 0,54 & 306 \\
        Juízes: convenções & Escore de convenções (camada 2) & 0,73 & 305 \\
        Juízes: nota geral & Larastan: erros novos & $-0{,}07$ & 291 \\
        Escore de convenções (camada 2) & Testes ocultos aprovados & 0,37 & 308 \\
        \bottomrule
    \end{tabular}
    \fonteautor
\end{table}

\textbf{Regras de convenção.} Das dez regras heurísticas (Apêndice~C), as
mais violadas no conjunto foram C4 (escritas de domínio por classes de ação:
cumprida em 39 das 148 execuções em que se aplicou), C3 (autorização
\textit{server-side} de toda ação Livewire: 36 de 87), C2 (validação com
\texttt{\#[Validate]}: 53 de 106), C8 (\texttt{\$fillable} ou \texttt{\$guarded}
declarados: 84 de 151) e C1 (código de produção acompanhado de teste: 193 de
308). As regras de escape em Blade (C7), de ausência de camadas de serviço
(C6), de interface em Bulma (C9) e de migrações geradas e reversíveis (C5)
foram cumpridas em quase todas as execuções em que se aplicaram. No piloto,
as 25 violações apontadas foram lidas manualmente contra o código, e todas
eram reais (C1 19 de 19, C2 2 de 2, C10 4 de 4); uma correção na rodada
completa fez a regra C8 reconhecer a sintaxe de atributo \texttt{\#[Fillable]},
que antes era contada como violação.

%-------------------------------------------------------------------------------
\subsection{Discussão}
\label{sec:discussao}

\textbf{O \textit{harness} muda como o código é escrito, não se a tarefa é
resolvida.} Os dois conjuntos de resultados apontam na mesma direção. Pelo
oráculo oculto, nenhum nível alterou o \textit{pass@1} de nenhum modelo; pelos
juízes e pela camada estática, todos os níveis elevaram convenções, testes e
casos de teste escritos, sem tocar em requisitos, segurança ou escopo. A
dissociação é coerente com o que o \textit{harness} contém: o \texttt{AGENTS.md}
não ensina o domínio nem corrige a leitura de um enunciado, mas prescreve como
o código deve ser organizado e exige que cada mudança venha com testes. O
efeito sobre a dimensão de testes é, provavelmente, o mais direto: a regra
``nenhum código de produção sem teste'' é lida e seguida pelos modelos de
capacidade intermediária e alta. O exemplo da T07 mostra o limite do oráculo:
um PR sem nenhum teste e um PR com quinze testes passam igualmente nos testes
ocultos, mas um revisor humano não os aprovaria da mesma forma.

\textbf{Escala explica o \textit{pass@1}; o \textit{harness} explica parte da
qualidade.} Entre modelos, o \textit{pass@1} vai de 11\% a 100\%; entre níveis,
para um mesmo modelo, a variação é de no máximo 7 pontos percentuais. A escolha
do modelo explica quase toda a variação em acerto, e o \textit{harness} não
compensa a diferença: o Haiku 5.5 não alcança o Sonnet 5.5 em nenhum nível, e o
gpt-oss-20b não sai do patamar de 1 a 2 tarefas em 14. Na qualidade, o quadro é
diferente. O DeepSeek v4.1 flash com \textit{harness} recebe dos juízes uma nota
geral (3,93 a 4,11) comparável à do Sonnet 5.5 sem \textit{harness} (3,96) e
escreve um número de testes (28 a 32 por execução) semelhante ao do Sonnet
com \textit{harness} (34). Nesse sentido estrito, e apenas nessa métrica, a
estrutura de apoio aproximou um modelo intermediário de um modelo maior. A
leitura mais defensável do lema \textit{scaffolding over scale}, à luz destes
dados, é portanto condicional: a estrutura de apoio não substitui a escala na
capacidade de resolver a tarefa, mas pode substituir parte da escala na
disciplina com que o código é entregue, desde que o modelo seja capaz de ler
e seguir as regras.

\textbf{O efeito concentra-se no meio da distribuição.} Quem mais ganha são
os modelos que partem de uma posição intermediária (DeepSeek e Muse). Os
modelos de ponta já escrevem testes e seguem convenções sem instrução, e o
\textit{harness} acrescenta pouco. Os modelos mais fracos (gpt-oss-20b e Beam)
não respondem: as regras estão na árvore, mas o modelo não as lê ou não as
segue, e a nota de testes permanece em 1,0. Há, assim, um limiar de capacidade
abaixo do qual a estrutura de apoio é inerte, o que contraria a versão mais
forte da hipótese, segundo a qual o \textit{harness} compensaria justamente os
modelos menores.

\textbf{Qual camada carrega o efeito.} A maior parte do ganho de qualidade e
do acréscimo de custo aparece já no nível \texttt{agents}. O nível
\texttt{skills} acrescenta um pouco em convenções e testes (os maiores $\Delta$
estão nele), o que é coerente com o conteúdo das \textit{skills}, que detalham
exatamente essas práticas. O nível \texttt{full} não acrescenta nada mensurável
e, em alguns casos, devolve parte do ganho; a configuração de LSP, cuja intenção
é reduzir leituras, não reduziu o consumo de \textit{tokens} de entrada em
relação ao \texttt{skills} de forma consistente. Com uma repetição, as
diferenças entre os três níveis com \textit{harness} estão dentro do ruído, e a
conclusão segura é apenas que a primeira camada responde pela maior parte do
efeito.

\textbf{Custo contra benefício.} O \textit{harness} encarece cada execução em
dezenas de pontos percentuais em turnos e \textit{tokens} de entrada, para um
\textit{pass@1} idêntico. Se o critério for apenas resolver a tarefa medida por
testes, o \textit{harness} é custo sem retorno nesta amostra. Se o critério for
entregar um PR com testes e dentro das convenções do projeto, o acréscimo
compra um ganho de 0,3 ponto na nota geral e de 0,5 a 0,8 ponto em testes, mais
dez casos de teste por mudança, o que, em um fluxo real, reduziria o trabalho
de revisão e de complementação posterior. Parte do custo adicional é, aliás, o
custo de escrever e executar os testes que o \texttt{AGENTS.md} exige. A
relação custo-benefício depende, portanto, do que se pede ao agente e de quem
revisa o resultado.

\textbf{Relação com a literatura.} O resultado dialoga com o achado de
\citeonline{yang2024sweagent} de que a interface do agente altera a taxa de
resolução: aqui, a intervenção foi textual e versionada no repositório, e não na
ferramenta, e o seu efeito apareceu na qualidade, não na resolução, em um
conjunto de tarefas em que a resolução já estava no teto. A dissociação entre
correção funcional e qualidade ecoa \citeonline{liu2023evalplus} e
\citeonline{pearce2022asleep}, que mostram, por caminhos diferentes, que passar
nos testes não esgota o que se espera de um código. Por fim, a concordância
entre juízes e a sua correlação com medidas independentes estão em linha com
o que \citeonline{zheng2023judging} e \citeonline{zhuo2024icescore} relatam, e o
cuidado com a autopreferência segue \citeonline{panickssery2024self}.

%-------------------------------------------------------------------------------
\subsection{Ameaças à validade}
\label{sec:ameacas}

\textbf{Validade de conclusão.} Com uma repetição por célula e 14 tarefas por
modelo, o desenho detecta apenas efeitos grandes sobre o \textit{pass@1}: o
teste de McNemar exige ao menos seis discordâncias na mesma direção para
$p < 0{,}05$, de modo que efeitos de 5 a 15 pontos percentuais passariam
despercebidos. Ausência de evidência não é evidência de ausência. A variância
de um mesmo modelo na mesma célula não foi estimada; as 20 células de
repetições 2 e 3 da rodada completa cancelada ficam no manifesto como amostra
pequena dessa variação. Na análise de qualidade foram feitas 39 comparações sem
correção para múltiplos testes; com um controle de taxa de falsas descobertas
\cite{benjamini1995fdr}, os efeitos em convenções, testes, casos de teste e nota
geral ($p \le 0{,}001$) resistiriam, mas os de Larastan, complexidade e
legibilidade (0,015 a 0,057) devem ser lidos como sugestivos. As cinco
subtarefas dos dois PRs compostos são correlacionadas (T04b e T04c incluem os
testes de T04a; T03a e T03b compartilham rotas), o que reduz o número efetivo
de pares.

\textbf{Validade interna.} O efeito de teto é a principal ameaça: quatro
modelos resolvem 93\% ou mais das tarefas sem \textit{harness}, e os testes
ocultos pareceram pouco discriminativos para modelos fortes (fração aprovada de
0,993 ou mais mesmo nas falhas). Tarefas mais difíceis seriam necessárias para
medir ganho de \textit{pass@1} acima do teto. O Beam ficou incompleto (29 de 56
células, de dois \textit{endpoints} do mesmo modelo) e o Qwen foi retirado, de
modo que esses modelos não entram nas comparações. O limite de memória de 6~GB
foi introduzido durante a rodada; nenhuma execução posterior morreu por
memória, e a anterior afetada foi refeita. Três reinícios do orquestrador
recomeçaram agentes em andamento como novas tentativas, sem duplicar células.
O modelo local rodou quantizado e com o analisador de formato do llama.cpp;
24 das suas 56 falhas são de formato de saída, uma limitação do conjunto
modelo mais servidor e não apenas do modelo.

\textbf{Validade de construto.} O oráculo é parcialmente escrito com IA:
cerca de 15\% dos \textit{commits} do repositório-sujeito são rastreáveis a
ferramentas de IA e os testes ocultos foram, às vezes, escritos com
assistência. Os testes ocultos acoplam-se a nomes: uma implementação correta
com nomes diferentes falha, o que o contrato do enunciado mitiga sem eliminar.
Os juízes são LLMs, com os vieses conhecidos; a concordância entre eles é alta,
mas eles podem concordar no mesmo erro, e não houve calibração com revisores
humanos. O Muse Spark ajudou a escrever o código do repositório-sujeito, treina
com os dados enviados e é, ao mesmo tempo, modelo avaliado e juiz de outros
modelos; a checagem de modelo confirmou que apenas o modelo sob teste respondeu
em cada execução, e o Muse nunca julgou o próprio código, mas uma afinidade de
estilo com o repositório não pode ser descartada. A divergência entre Larastan
e juízes (correlação $-0{,}53$ no piloto e $-0{,}07$ na rodada completa) indica
que a contagem de erros novos do Larastan captura algo que os juízes não veem,
ou vice-versa; o Larastan não está instalado no projeto e os agentes não o
executam, e a aplicação não inicializa em parte das execuções do gpt-oss, o que
reduz o $n$ dessa métrica. As regras C1 a C10 são heurísticas: a validação
manual no piloto encontrou apenas violações reais, mas regras que não se
aplicaram no piloto (C5 e C8) só foram exercitadas na rodada completa, e uma
delas precisou de correção.

\textbf{Validade externa.} As tarefas vêm de um único repositório privado, de
um único autor, em uma única pilha (Laravel, Livewire, Bulma), e o
\textit{harness} é o desse repositório. Nada aqui autoriza generalizar para
outros repositórios, linguagens ou \textit{harnesses}. Os \textit{runners} são
diferentes entre famílias de modelos, de modo que comparações entre modelos
misturam modelo e \textit{runner}; apenas as comparações de níveis dentro de
cada modelo são controladas. O \textit{sandbox} tinha internet liberada; a
natureza privada do repositório e o detector cobrem o risco de vazamento, e
nenhum vazamento foi observado. O preâmbulo não interativo é igual para todos os
modelos e níveis, mas é uma condição específica deste experimento.

% ==============================================================================
% CAPÍTULO 5 - CONSIDERAÇÕES FINAIS
% ==============================================================================

\section{Considerações Finais}
\label{sec:consideracoes}

Este trabalho perguntou se um \textit{harness} em nível de repositório compensa
a escala do modelo em agentes de código. Para respondê-la, construiu um
\textit{benchmark} de ablação com 14 tarefas derivadas de \textit{pull requests}
reais de uma aplicação Laravel privada, avaliadas por testes ocultos; um
orquestrador que executou seis modelos em quatro níveis cumulativos de
\textit{harness}, em \textit{sandbox} e com registro reprodutível; e uma análise
de qualidade de código em três camadas sobre 309 execuções e 836 revisões de
juízes. Os cinco objetivos específicos foram atingidos: o \textit{benchmark}, os
níveis e a execução foram realizados como planejado, com a redução de três
repetições para uma registrada e justificada; a análise de qualidade e a
análise estatística pareada foram conduzidas e seus resultados discutidos.

A resposta à pergunta é dupla e depende da métrica. Na taxa de resolução, o
\textit{harness} não compensou a escala: quatro dos seis modelos já resolviam
95\% ou mais das tarefas sem nenhuma estrutura de apoio, nenhum nível alterou o
\textit{pass@1} de nenhum modelo (McNemar exato, $p = 1{,}0$ em 70 pares por
nível) e cada execução ficou mais cara em turnos e \textit{tokens}. Entre
modelos, o \textit{pass@1} variou de 11\% a 100\%; entre níveis de um mesmo
modelo, no máximo 7 pontos percentuais. Na qualidade do código, o
\textit{harness} teve efeito consistente: a nota geral dos juízes subiu 0,31 a
0,34 ponto, a nota de convenções 0,46 a 0,52 e a de testes 0,60 a 0,87, com
cerca de dez casos de teste a mais por execução, menos erros novos de análise
estática e menor complexidade, e sem alteração em requisitos, segurança ou
escopo. O efeito concentrou-se nos modelos de capacidade intermediária
(DeepSeek v4.1 flash e Muse Spark 1.3), foi pequeno nos modelos de ponta, que
já partem alto, e nulo nos modelos mais fracos (gpt-oss-20b e Beam), que não
seguem as regras. A primeira camada, as regras sempre ativas, respondeu pela
maior parte do efeito e do custo; as \textit{skills} acrescentaram pouco e a
configuração de LSP nada mensurável. Em síntese, nesta amostra a escala
explica se a tarefa é resolvida e o \textit{harness} explica parte de como o
código é entregue.

As contribuições são o \textit{benchmark} e seus enunciados, o orquestrador, o
pacote de análise de qualidade com a rubrica e as regras de convenção, e os
dados de todas as execuções e revisões, disponíveis publicamente, com exceção
dos \textit{commits} do repositório-sujeito, que é privado. Do ponto de vista
metodológico, o trabalho mostra que um oráculo de testes, sozinho, não
distingue entregas que um revisor distinguiria, e que juízes LLM de três
famílias, com evidência obrigatória, controle de autopreferência e validação
cruzada por medidas objetivas, podem preencher essa lacuna com concordância
substancial.

As limitações são as de um experimento de uma repetição, 14 tarefas de um
único repositório privado e efeito de teto: o desenho só detectaria efeitos
grandes sobre o \textit{pass@1}, e nada aqui se generaliza para outros
repositórios, linguagens ou \textit{harnesses}. O oráculo foi parcialmente
escrito com IA, os juízes são LLMs sem calibração humana, as regras de
convenção são heurísticas, o Beam ficou incompleto e as comparações múltiplas
não foram corrigidas.

Como trabalhos futuros, propõem-se: executar mais repetições por célula, para
estimar a variância de um mesmo modelo e ganhar poder estatístico; selecionar
tarefas mais difíceis, fora do teto, nas quais um ganho de \textit{pass@1} seja
mensurável; repetir o protocolo em outros repositórios e pilhas, com seus
próprios \textit{harnesses}, para testar a generalidade dos achados; fazer a
ablação de partes do \texttt{AGENTS.md} (apenas a regra de testes, apenas as
convenções, apenas a política de segurança), para identificar quais frases
carregam o efeito; e calibrar os juízes LLM com revisores humanos em uma
amostra das mudanças, para medir a validade das notas diretamente. Uma rodada
própria para os modelos locais, com o Qwen retirado desta, e a conclusão da
matriz do Beam completariam o quadro dos modelos de menor capacidade, onde a
hipótese \textit{scaffolding over scale} tem mais a provar.


% ==========================================================================
% ELEMENTOS PÓS-TEXTUAIS
% ==========================================================================
\postextual

% --- Declaração sobre o uso de IA generativa (opcional) ---
% Remova ou comente a linha abaixo caso nenhuma ferramenta tenha sido utilizada.
% ==============================================================================
% DECLARAÇÃO SOBRE O USO DE IA GENERATIVA E TECNOLOGIAS ASSISTIDAS POR IA
% ==============================================================================

\section*{Declaração sobre o uso de IA generativa}
\label{sec:declaracao-uso-ia}

Ferramentas de inteligência artificial generativa foram utilizadas neste
trabalho de quatro formas distintas, descritas também na
\autoref{sec:uso-ia}.

Primeiro, modelos de linguagem foram o \textbf{objeto} do experimento: Claude
Sonnet 5.5 e Claude Haiku 5.5 (Anthropic, via Claude Code 2.1.312), Muse Spark
1.3, DeepSeek v4.1 flash e Beam (via OpenCode v2.2.0) e gpt-oss-20b (OpenAI,
executado localmente com llama.cpp) escreveram o código avaliado.

Segundo, a ferramenta Claude Code, desenvolvida pela Anthropic, com o modelo
Claude Opus 5.5, foi utilizada na construção do orquestrador, dos
\textit{scripts} de análise e dos relatórios intermediários do
\textit{benchmark}, sob direção e revisão do autor. O Muse Spark foi utilizado
como ferramenta na criação de parte do código do repositório-sujeito, anterior
a este trabalho.

Terceiro, os modelos Claude Opus 5.5 (Anthropic), Muse Spark 1.3 e DeepSeek
v4.1 flash atuaram como \textbf{juízes} na terceira camada da análise de
qualidade de código, inclusive sobre o código da própria família, como controle
de autopreferência, sob as salvaguardas descritas na
\autoref{sec:qualidade-metodo}.

Quarto, o texto deste artigo foi redigido com o auxílio do modelo Claude Fable
5.1, desenvolvido pela Anthropic, a partir dos relatórios, dos dados e da
documentação produzidos durante a pesquisa, e foi integralmente revisado pelo
autor. Todo o conteúdo produzido com o auxílio dessas ferramentas foi
analisado, revisado e validado pelo autor, que assume integral responsabilidade
pelo conteúdo desta publicação.


% --- Referências bibliográficas ---
\bibliography{referencias}
\newpage


% --- Anexos ---
% Este trabalho não possui anexos (documentos de terceiros).

% --- Apêndices ---
\begin{apendicesenv}
    % %Imprime uma página indicando o início dos apêndices
    % \partapendices

    % ==============================================================================
% APÊNDICE A
% ==============================================================================

\section{As 14 tarefas do \textit{benchmark}}
\label{apd:tarefas}

O \autoref{qdr:tarefas} lista as tarefas da rodada reduzida. O porte segue o
número de arquivos de produção da solução de referência (pequena até 10, média
de 11 a 25, grande acima de 25). A coluna Arq./lin. traz o tamanho da solução de
referência (arquivos alterados e linhas inseridas, incluindo testes). As tarefas T03a e
T03b são subtarefas do PR de autenticação e T04a, T04b e T04c são subtarefas do
PR de entidades; os PRs compostos não entram no conjunto. Duas tarefas (T10 e
T11) vêm de \textit{branches} não mescladas.

\begin{quadro}[!htb]
    \centering
    \ABNTEXfontereduzida
    \caption{\label{qdr:tarefas}Tarefas do \textit{benchmark}, com categoria, porte, origem, tamanho da referência e observações}
    \begin{tabular}{|p{0.7cm}|p{4.1cm}|p{1.1cm}|p{1.2cm}|p{1.3cm}|p{1.1cm}|p{4.4cm}|}
        \hline
        \textbf{Id} & \textbf{Título} & \textbf{Categ.} & \textbf{Porte} & \textbf{Origem} & \textbf{Arq./lin.} & \textbf{Observações} \\
        \hline
        T01 & Notificador de alertas no Discord e canais de \textit{log} & ops & pequena & PR 17 & 10/400 & Bem especificada; 2 testes ocultos novos. \\
        \hline
        T02 & Menção a um usuário configurado em todo alerta & ops & pequena & PR 30 & 6/103 & Muito pequena; modifica um teste existente; item de calibração. \\
        \hline
        T03a & Autenticação de sessão (login, registro, redefinição, verificação) & auth & grande & PR 19 (sub.) & 31/1\,062 & Subtarefa de T03; 3 testes novos; acoplamento a nomes de rota. \\
        \hline
        T03b & Redefinição de senha e e-mail transacional com tema & auth & pequena & PR 19 (sub.) & 11/584 & Subtarefa de T03; 1 teste novo; chaves de tradução. \\
        \hline
        T04a & Entidades, associações com papéis e \textit{policies} & tenancy & média & PR 31 (sub.) & 15/768 & Subtarefa de T04; 4 testes novos; coluna gerada no banco. \\
        \hline
        T04b & Entidades mais \textit{onboarding}, sessão de locatário e seletor & tenancy & média & PR 31 (sub.) & 33/1\,566 & Inclui os testes de T04a; acoplamento de marcação. \\
        \hline
        T04c & Entidades mais painel, gestão de membros e configurações & tenancy & média & PR 31 (sub.) & 28/1\,831 & Inclui os testes de T04a; testes de adulteração em Livewire. \\
        \hline
        T05 & Associação por convite com e-mail e modal de membro & tenancy & média & PR 34 & 24/1\,456 & Modifica um teste existente; exige remover comportamento antigo. \\
        \hline
        T06 & Filas, senhas e ação de emissão de senha & domain & média & PR 33 & 14/679 & 14 testes de domínio, sem UI; sequência de códigos e escrita única. \\
        \hline
        T07 & Chamada da próxima senha com modos e justiça de 3 chamadas & domain & pequena & PR 37 & 6/390 & 15 testes; algorítmica; desempate por id. \\
        \hline
        T08 & Ações de gestão de filas, permissões por fila e \textit{policy} & domain & média & PR 39 & 17/1\,149 & Modifica um teste existente ao exigir autorização. \\
        \hline
        T09 & Telas Livewire de gestão de filas e painel do atendente & UI & média & PR 40 & 15/2\,013 & Grande superfície de UI; acoplamento de marcação; 2 testes modificados. \\
        \hline
        T10 & API BFF do totem (pareamento, JWT, filas, prioridades, emissão) & API & grande & \textit{branch} 42$^{a}$ & 38/1\,468 & Dependência JWT; CORS; códigos de erro exatos. \\
        \hline
        T11 & Prefixos de código por prioridade com sequências diárias & domain & pequena & \textit{branch} 46$^{a}$ & 14/518 & Modifica 2 testes existentes; muitos casos de borda. \\
        \hline
    \end{tabular}
    \legend{Fonte: elaborado pelo autor (2026). $^{a}$ \textit{Branch} não mesclada (tarefa \textit{held-out}).}
\end{quadro}

    \clearpage
    % ==============================================================================
% APÊNDICE B
% ==============================================================================

\section{Rubrica dos juízes}
\label{apd:rubrica}

Cada juiz recebe o enunciado, as seções normativas do \texttt{AGENTS.md}, o
\textit{diff} da mudança, a rubrica e um esquema JSON de saída. Ele atribui a
cada dimensão uma nota de 1 a 5, com evidência obrigatória
(\texttt{caminho:linha} ou trecho citado), e uma nota geral que responde à
pergunta ``você aprovaria este \textit{pull request}?'' (5: aprova como está;
4: aprova com observações; 3: precisa de mudanças; 2: precisa de retrabalho
substancial; 1: rejeita). O juiz também lista os problemas encontrados, com
severidade, dimensão, local e explicação, e declara a confiança na própria
revisão. O \autoref{qdr:rubrica} resume as âncoras; as notas 2 e 4 são usadas
para casos intermediários. O juiz é instruído a não recompensar o comprimento,
a não comparar com outras mudanças e a não tentar adivinhar quem escreveu o
código.

\begin{quadro}[htb]
    \centering
    \ABNTEXfontereduzida
    \caption{\label{qdr:rubrica}Dimensões da rubrica e âncoras das notas 5, 3 e 1}
    \begin{tabular}{|p{2.3cm}|p{4.2cm}|p{3.4cm}|p{4.0cm}|}
        \hline
        \textbf{Dimensão} & \textbf{5} & \textbf{3} & \textbf{1} \\
        \hline
        Requisitos & Todos os requisitos e o contrato atendidos; casos de borda tratados & Caminho principal feito; algum requisito, caso de borda ou detalhe do contrato ausente ou errado & Maioria dos requisitos ausente, contrato errado ou código que não funciona \\
        \hline
        Convenções & Segue toda convenção aplicável (ações para escritas de domínio, \texttt{\#[Validate]}, \textit{policies}, migrações geradas, \texttt{\$fillable}, componentes Blade, Bulma, docs em sincronia) & Segue a maioria; um ou dois desvios claros & Ignora as convenções (camadas de serviço, lógica em \textit{views}, padrões refeitos à mão) \\
        \hline
        Segurança & Nenhuma fraqueza; autorização e validação onde necessárias & Uma lacuna de impacto limitado & Falha explorável (escrita não autorizada, saída sem escape, confiança em ids do cliente) \\
        \hline
        Testes & Testes cobrem todos os requisitos e os caminhos de falha e autorização & Alguns testes, só caminho feliz, ou casos importantes ausentes & Nenhum teste para o novo comportamento, ou testes enfraquecidos ou apagados \\
        \hline
        Simplicidade & Enxuto e direto; reutiliza o que existe & Alguma duplicação ou indireção desnecessária & Sobre-engenharia pesada ou cópia e cola \\
        \hline
        Legibilidade & Fácil de ler e de mudar & Compreensível com esforço; métodos longos ou nomes pouco claros & Difícil de seguir; código morto, saída de depuração, estilo inconsistente \\
        \hline
        Escopo & Tudo no \textit{diff} serve à tarefa & Algumas edições não relacionadas & Grandes mudanças ou funcionalidades não pedidas \\
        \hline
    \end{tabular}
    \fonteautor
\end{quadro}

    \clearpage
    % ==============================================================================
% APÊNDICE C
% ==============================================================================

\section{Regras de convenção C1 a C10}
\label{apd:regras}

As dez regras da camada 2 foram derivadas do \texttt{AGENTS.md} e verificadas
automaticamente no \textit{patch} de cada execução. Cada regra só se aplica
quando a mudança toca o seu assunto (por exemplo, C2 e C3 só se aplicam quando
há um componente Livewire no \textit{diff}), e o escore da execução é a fração
das regras aplicáveis cumpridas. Toda violação guarda a evidência que a
motivou. A \autoref{tab:regras} lista as regras e, para as 309 execuções da
rodada, quantas vezes cada uma se aplicou e foi cumprida.

\begin{table}[htb]
    \centering
    \ABNTEXfontereduzida
    \caption{Regras de convenção, número de execuções em que cada uma se aplicou e em que foi cumprida (309 execuções)}
    \label{tab:regras}
    \begin{tabular}{lp{9.2cm}rr}
        \toprule
        \textbf{Regra} & \textbf{Descrição} & \textbf{Aplica} & \textbf{Cumpre} \\
        \midrule
        C1 & Código de produção acompanhado de teste (\textit{test-first}) & 308 & 193 \\
        C2 & Componentes Livewire validam com \texttt{\#[Validate]} e \texttt{\$this->validate()} & 106 & 53 \\
        C3 & Toda ação pública de Livewire é autorizada no servidor & 87 & 36 \\
        C4 & Escritas de domínio passam por classes em \texttt{app/Actions} & 148 & 39 \\
        C5 & Migrações geradas (nome com \textit{timestamp}) e reversíveis (\texttt{down()}) & 171 & 168 \\
        C6 & Nenhuma camada de serviço ou repositório & 261 & 260 \\
        C7 & Blade escapa a saída (\texttt{\{\{ \}\}}, nunca \texttt{\{!! !!\}}) & 130 & 130 \\
        C8 & Modelos novos declaram \texttt{\$fillable} ou \texttt{\$guarded} & 151 & 84 \\
        C9 & Interface em Bulma, sem classes utilitárias de CSS nem bibliotecas de UI em JavaScript & 130 & 129 \\
        C10 & Autorização por \textit{policies} ou \textit{gates}, não por verificação manual de papel & 149 & 106 \\
        \bottomrule
    \end{tabular}
    \fonteautor
\end{table}

    \clearpage
    % ==============================================================================
% APÊNDICE D
% ==============================================================================

\section{Dados complementares da rodada reduzida}
\label{apd:complementares}

Este apêndice reúne as medianas absolutas de custo por execução, que a
\autoref{fig:custo} apresenta de forma relativa, e a linha do tempo dos
incidentes e decisões da rodada reduzida, mencionada na
\autoref{sec:orquestrador}.

\begin{table}[!htb]
    \centering
    \ABNTEXfontereduzida
    \caption{Custo por execução (medianas) por modelo e nível}
    \label{tab:custo}
    \begin{tabular}{llrrrr}
        \toprule
        \textbf{Modelo} & \textbf{Nível} & \textbf{Turnos} & \textbf{Tokens de saída} & \textbf{Tokens de entrada} & \textbf{Tempo (min)} \\
        \midrule
        Sonnet 5.5 & \texttt{none} & 18 & 22\,502 & 905\,810 & 3,1 \\
         & \texttt{agents} & 30 & 30\,112 & 1\,468\,970 & 4,1 \\
         & \texttt{skills} & 29 & 31\,398 & 1\,507\,708 & 4,5 \\
         & \texttt{full} & 26 & 27\,909 & 1\,363\,034 & 4,6 \\
        \midrule
        Haiku 5.5 & \texttt{none} & 62 & 73\,469 & 3\,767\,150 & 5,9 \\
         & \texttt{agents} & 70 & 86\,608 & 5\,397\,860 & 8,3 \\
         & \texttt{skills} & 74 & 85\,562 & 6\,428\,217 & 9,0 \\
         & \texttt{full} & 71 & 87\,200 & 5\,525\,890 & 8,1 \\
        \midrule
        Muse Spark 1.3 & \texttt{none} & 62 & 20\,360 & 4\,125\,950 & 4,1 \\
         & \texttt{agents} & 64 & 18\,510 & 4\,354\,079 & 8,6 \\
         & \texttt{skills} & 74 & 18\,452 & 4\,700\,618 & 5,1 \\
         & \texttt{full} & 67 & 19\,244 & 4\,989\,247 & 5,8 \\
        \midrule
        DeepSeek v4.1 flash & \texttt{none} & 58 & 17\,518 & 4\,256\,746 & 9,3 \\
         & \texttt{agents} & 68 & 20\,720 & 4\,962\,629 & 8,7 \\
         & \texttt{skills} & 67 & 20\,570 & 6\,097\,903 & 8,3 \\
         & \texttt{full} & 76 & 22\,686 & 6\,913\,314 & 8,4 \\
        \midrule
        gpt-oss-20b & \texttt{none} & 42 & 10\,990 & 693\,686 & 2,5 \\
         & \texttt{agents} & 38 & 10\,620 & 691\,718 & 1,8 \\
         & \texttt{skills} & 46 & 13\,660 & 977\,582 & 2,4 \\
         & \texttt{full} & 46 & 11\,222 & 982\,754 & 2,3 \\
        \bottomrule
    \end{tabular}
    \legend{Fonte: elaborado pelo autor (2026). Tokens de entrada incluem leitura de \textit{cache}. O tempo é o de parede do agente e não é comparável entre modelos.}
\end{table}

\begin{quadro}[!htb]
    \centering
    \ABNTEXfontereduzida
    \caption{\label{qdr:incidentes}Incidentes da rodada reduzida (29 de outubro de 2026) e decisões tomadas}
    \begin{tabular}{|l|p{5.6cm}|p{6.9cm}|}
        \hline
        \textbf{Hora} & \textbf{Evento} & \textbf{Decisão} \\
        \hline
        00h37 & Detector de contaminação marca um \textit{link} para um PR escrito no \texttt{CHANGELOG.md} & \textit{Links} para o repositório-sujeito deixam de contar; o veredito passa a ser recalculado do \textit{log} bruto completo; 2 falsos positivos voltam aos dados. \\
        \hline
        00h44 e 02h54 & DeepSeek tenta ler o repositório-sujeito no GitHub (arquivo bruto; depois a API) & Repositório privado, sem credencial no \textit{sandbox}: 404, sem vazamento. Tentativas contadas à parte como comportamento observado. \\
        \hline
        01h24 & Contêiner de um agente chega a 23~GB de \textit{swap} e é morto pelo sistema & Limite de 6~GB de memória por contêiner; a célula é refeita como infraestrutura. \\
        \hline
        01h52 & Qwen 3.6 35B-A3B leva cerca de 23 minutos por célula na GPU única & Modelo retirado da rodada (restariam cerca de 20 horas); 3 células ficam só como nota. \\
        \hline
        02h35 & Provedor do Beam fora do ar: agentes mudos por 10 a 15 minutos até o \textit{timeout} & \textit{Watchdog} de silêncio (8 minutos) trata o travamento como limite de taxa; 6 falsos \textit{timeouts} reenfileirados, nunca contados como dado. \\
        \hline
        03h22 & CLI do OpenCode sai com erro depois de o agente terminar & Classificado como infraestrutura e refeito; a decisão não dependeu do resultado. \\
        \hline
        07h08 a 07h45 & Beam retomado por outro \textit{endpoint} do mesmo modelo; cota gratuita esgota após 6 células & Beam encerrado com 29 das 56 células; apresentado só de forma descritiva. \\
        \hline
    \end{tabular}
    \fonteautor
\end{quadro}

\end{apendicesenv}

% --- Índice remissivo (opcional) ---
% \phantompart
% \printindex

\end{document}
